\documentclass[12pt,a4paper,UTF8]{ctexart}
\usepackage[a4paper,top=22mm,bottom=22mm,left=22mm,right=22mm,headheight=15pt]{geometry}
\newcommand{\Needspace}[1]{\par\begingroup\skip0=\lastskip\vskip-\skip0\dimen0=#1\relax\advance\dimen0 by\skip0\vskip\dimen0\penalty-100\vskip-\dimen0\vskip\skip0\penalty9999\endgroup}
\usepackage{amsmath,amssymb,mathtools,bm}
\usepackage{graphicx,xcolor,booktabs,longtable,array,enumitem,fancyhdr,caption,float,microtype,setspace}
\usepackage[hidelinks,unicode]{hyperref}
\usepackage[most]{tcolorbox}
\definecolor{ContestBlue}{HTML}{17365D}
\hypersetup{pdftitle={第078套},pdfauthor={}}
\setstretch{1.24}
\setlength{\parindent}{2em}
\setlength{\parskip}{0.35em}
\setlength{\emergencystretch}{2em}
\setlist{itemsep=0.35em,topsep=0.4em}
\captionsetup{font=small,labelfont=bf,justification=centering}
\setcounter{secnumdepth}{0}
\pagestyle{fancy}
\fancyhf{}
\fancyhead[L]{\small 第078套}
\fancyfoot[C]{\thepage}
\renewcommand{\headrulewidth}{0.4pt}
\ctexset{
  section={format=\Large\bfseries\color{ContestBlue},beforeskip=1.2em,afterskip=0.7em},
  subsection={format=\large\bfseries\color{ContestBlue},beforeskip=1.0em,afterskip=0.5em},
  subsubsection={format=\normalsize\bfseries,beforeskip=0.8em,afterskip=0.35em}
}
\usepackage{calc}
\newcounter{none}
\let\originalsubsection\subsection
\renewcommand{\subsection}{\Needspace{6\baselineskip}\originalsubsection}
\makeatletter
\newcommand{\tightlist}{\setlength{\itemsep}{0pt}\setlength{\parskip}{0pt}}
\newcommand{\maxwidth}{\ifdim\Gin@nat@width>\linewidth\linewidth\else\Gin@nat@width\fi}
\newcommand{\maxheight}{\ifdim\Gin@nat@height>0.72\textheight 0.72\textheight\else\Gin@nat@height\fi}
\newcommand{\pandocbounded}[1]{\begingroup\setkeys{Gin}{width=\maxwidth,height=\maxheight,keepaspectratio}#1\endgroup}
\makeatother
\begin{document}
\begin{titlepage}
  \centering
  \vspace*{58mm}
  {\Huge\bfseries\color{ContestBlue} 第078套\par}
  \vfill
  \vspace*{28mm}
\end{titlepage}
\section{第 1 题}\label{ux7b2c-1-ux9898}

无限大接地导体平面上有一个半径为 \(R\)
的半球形突起。以半球球心为原点建立直角坐标系 \(Oxyz\)，导体平面为
\(z=0\)。在

\[P(\sqrt2R,0,\sqrt2R)\]

处放置一点电荷 \(Q\)。求该电荷所受静电力在 \(x,y,z\) 三个方向上的分量。

\begin{figure}[H]
\centering
\includegraphics[width=0.72\linewidth,height=\textheight,keepaspectratio,alt={图 1-1　接地平面、半球形突起与点电荷的位置关系。}]{../figures/question-01.pdf}
\caption*{图 1-1　接地平面、半球形突起与点电荷的位置关系。}
\end{figure}

\subsection{一、题型定位}\label{ux4e00ux9898ux578bux5b9aux4f4d}

平面镜像和球面反演要配成同一组电荷，才能同时满足两段接地边界。先逐对检验电势，再把三个像电荷的力分量相加。

\subsection{二、知识点}\label{ux4e8cux77e5ux8bc6ux70b9}

\begin{itemize}
\tightlist
\item
  接地平面的平面镜像。
\item
  接地球面的 Kelvin 镜像。
\item
  唯一性定理。
\item
  库仑力矢量分解。
\item
  对称性对力分量的限制。
\end{itemize}

\subsection{三、物理图像}\label{ux4e09ux7269ux7406ux56feux50cf}

半球面与平面相交于圆周，两部分属于同一接地导体。真实电荷位于 \(xOz\)
平面内，因此整个边界条件关于 \(xOz\) 平面对称，必有
\(F_y=0\)。其余两个分量则由三个像电荷共同决定。

\subsection{四、解题入口}\label{ux56dbux89e3ux9898ux5165ux53e3}

先关于平面作镜像，把半空间边界处理掉；再分别对真实电荷和平面像电荷关于半径为
\(R\) 的球面作 Kelvin 镜像。最后逐个计算三个像电荷对真实电荷的作用力。

\subsection{五、解析}\label{ux4e94ux89e3ux6790}

\begin{figure}[H]
\centering
\includegraphics[width=0.718\linewidth,height=\textheight,keepaspectratio,alt={图 1-2　同时满足平面与半球面零电势边界的三个像电荷。}]{../figures/solution-01.pdf}
\caption*{图 1-2　同时满足平面与半球面零电势边界的三个像电荷。}
\end{figure}

图 1-2 以球心 \(O\) 为共同原点，标出真实电荷 \(Q\)、平面像
\(q_1=-Q\)、真实电荷的 Kelvin 像 \(q_2=-Q/2\)，以及 \(q_1\) 的 Kelvin 像
\(q_3=+Q/2\)。下面计算力时，各位置矢量以及从像电荷指向真实电荷的位移矢量均严格沿用图中的方向。

设

\[a=\sqrt2R,\]

真实电荷的位置矢量为

\[\bm r_P=(a,0,a).\]

它到球心的距离为

\[r_P=\sqrt{a^2+a^2}=2R.\]

\Needspace{3\baselineskip}
\textbf{1. 平面镜像}

关于 \(z=0\) 平面镜像，得到

\[q_1=-Q,\qquad \bm r_1=(a,0,-a).\]

真实电荷与 \(q_1\) 的间距为 \(2a=2\sqrt2R\)，方向沿 \(-z\)。所以

\[
\bm F_1
=-\frac{1}{4\pi\varepsilon_0}\frac{Q^2}{(2\sqrt2R)^2}\bm e_z
=-\frac{Q^2}{32\pi\varepsilon_0R^2}\bm e_z.
\]

\Needspace{3\baselineskip}
\textbf{2. 真实电荷的球面镜像}

接地球面的镜像公式是

\[q'=-\frac{R}{r_P}Q,\qquad \bm r'=\frac{R^2}{r_P^2}\bm r_P.\]

由于 \(r_P=2R\)，有

\[q_2=-\frac Q2,\]

\[
\bm r_2=\frac14(a,0,a)
=\left(\frac{\sqrt2R}{4},0,\frac{\sqrt2R}{4}\right).
\]

从 \(q_2\) 指向真实电荷的位移为

\[
\bm d_2=\bm r_P-\bm r_2
=\frac{3a}{4}(1,0,1).
\]

其长度为

\[d_2=\frac{3a}{4}\sqrt2=\frac{3R}{2}.\]

所以

\[
\begin{aligned}
\bm F_2
&=\frac{1}{4\pi\varepsilon_0}
\frac{Qq_2}{d_2^3}\bm d_2\\
&=-\frac{1}{4\pi\varepsilon_0}
\frac{\sqrt2Q^2}{9R^2}(1,0,1).
\end{aligned}
\]

\Needspace{3\baselineskip}
\textbf{3. 平面像电荷的球面镜像}

\(q_1=-Q\) 到球心的距离同样为 \(2R\)，故它的球面镜像为

\[q_3=-\frac{R}{2R}q_1=+\frac Q2,\]

\[
\bm r_3=\left(\frac{\sqrt2R}{4},0,-\frac{\sqrt2R}{4}\right).
\]

位移矢量为

\[
\bm d_3=\bm r_P-\bm r_3
=\frac{\sqrt2R}{4}(3,0,5).
\]

于是

\[
d_3=\frac{\sqrt2R}{4}\sqrt{3^2+5^2}
=\frac{\sqrt{68}}{4}R
=\frac{\sqrt{17}}{2}R.
\]

代入库仑定律：

\[
\begin{aligned}
\bm F_3
&=\frac{1}{4\pi\varepsilon_0}
\frac{Q(Q/2)}{d_3^3}\bm d_3\\
&=\frac{1}{4\pi\varepsilon_0}
\frac{\sqrt{34}Q^2}{289R^2}(3,0,5).
\end{aligned}
\]

\Needspace{3\baselineskip}
\textbf{4. 合成三个分量}

把三项相加，得到

\[
\boxed{
F_x=\frac{Q^2}{4\pi\varepsilon_0R^2}
\left(-\frac{\sqrt2}{9}+\frac{3\sqrt{34}}{289}\right)
},
\]

\[\boxed{F_y=0},\]

\[
\boxed{
F_z=\frac{Q^2}{4\pi\varepsilon_0R^2}
\left(-\frac18-\frac{\sqrt2}{9}+\frac{5\sqrt{34}}{289}\right)
}.
\]

\Needspace{3\baselineskip}
\textbf{边界条件的直接检查}

在 \(z=0\) 平面上，\(Q\) 与 \(-Q\) 到任一点距离相等，\(-Q/2\) 与
\(+Q/2\) 到任一点距离也相等，两对电势分别抵消。在半径为 \(R\)
的球面上，每个外部电荷与自己的 Kelvin
像电荷构成零电势对。由静电场唯一性定理，这组像电荷给出的上半空间电场就是所求物理解。

\Needspace{3\baselineskip}
\textbf{矢量计算的分量表}

把公共因子 \(K=Q^2/(4\pi\varepsilon_0R^2)\)
提出，三个像电荷的贡献可整理为

\[
\begin{array}{c|ccc}
&F_x/K&F_y/K&F_z/K\\ \hline
q_1&0&0&-1/8\\
q_2&-\sqrt2/9&0&-\sqrt2/9\\
q_3&3\sqrt{34}/289&0&5\sqrt{34}/289
\end{array}
\]

这张表能直接检查：每一项都具有 \(K\) 的力的量纲，且所有 \(y\)
分量严格为零。

\Needspace{3\baselineskip}
\textbf{像电荷位置为什么都在导体内部}

球面像电荷的位置半径满足

\[r'=\frac{R^2}{r_P}.\]

本题 \(r_P=2R\)，所以两个球面像电荷到球心的距离都是 \(R/2\)，严格小于
\(R\)。平面像电荷位于 \(z<0\)
区域。它们全都落在求解区域之外，因此不会在真实空间额外引入点电荷奇点。

\Needspace{3\baselineskip}
\textbf{力的方向直觉}

接地导体对正电荷总体表现为吸引。结果中

\[F_x<0,\qquad F_z<0,\]

分别指向半球轴线和导体表面。数值上

\[
\frac{F_x}{K}\approx-0.0966,\qquad
\frac{F_z}{K}\approx-0.1814,
\]

方向与几何直觉一致。若算出 \(F_z>0\)，通常是把平面像电荷的符号写错。

\Needspace{3\baselineskip}
\textbf{从势能求力的交叉检验}

真实电荷在像电荷产生的电势中的静电能为

\[U=\frac12Q\varphi_{\rm image}(\bm r_P).\]

像电荷的电量和位置都会随真实电荷位置改变，因此用势能求导时必须保留这一依赖，并保留系数
\(1/2\)。对一般位置 \((a,0,b)\) 写出 \(U(a,b)\) 后，

\[F_x=-\frac{\partial U}{\partial a},\qquad
F_z=-\frac{\partial U}{\partial b},\]

再代入 \(a=b=\sqrt2R\)，可重新得到正文结果。这提供了独立的交叉核验。

\Needspace{3\baselineskip}
\textbf{近表面极限}

当电荷逐渐靠近平面且远离半球时，最近的平面像电荷主导，结果趋近

\[F_z\simeq-\frac{Q^2}{16\pi\varepsilon_0z^2}.\]

当电荷靠近半球面时，最近的 Kelvin
像电荷主导，局部曲面也趋于平面。两个极限均表现为指向导体的吸引力。

\Needspace{3\baselineskip}
\textbf{一般位置的能量核对}

令实电荷位于 \((a,0,b)\)、\(s=\sqrt{a^2+b^2}>R\)，两个球面像的反演系数为
\(R^2/s^2\)。感应势能可由

\[U=\frac{kQ^2}{2}\left[-\frac1{2b}-\frac{R}{s^2-R^2}
+\frac{R}{\sqrt{(s^2-R^2)^2+4R^2b^2}}\right]\]

求导。这里 \(k=1/(4\pi\varepsilon_0)\)。分别取
\(-\partial U/\partial a\) 和 \(-\partial U/\partial b\)，再代
\(a=b=\sqrt2R\)，得到正文两分量。图 1-2 的球面像距离为
\(R/2\)，负像的力指向左下，正像的力指向右上；两箭头不抵消，因为到真实电荷的距离不同。

\Needspace{3\baselineskip}
\textbf{两组配对分别检验球面和接地平面}

平面上任一点 \(z=0\)，真实 \(Q\) 与平面像 \(-Q\)
到该点距离相等，电势抵消；两个球面像 \(-Q/2,+Q/2\)
也关于平面对称，第二对同样抵消。

在半径 \(R\) 的球面上，实电荷距球心为 \(s\)、反演像位置为
\(R^2\bm r_P/s^2\)。对球面点 \(\bm r\)，有

\[\left|\bm r-\frac{R^2}{s^2}\bm r_P\right|
=\frac R s|\bm r-\bm r_P|.\]

像电荷大小为
\(-RQ/s\)，所以它的势恰好与真实电荷的势相反。平面像与它的球面反演像也成同样一对。因此四电荷势同时满足两段零边界，并在无穷远趋零，唯一性定理保证真实区域的解。

求实电荷受力时仅叠加三个像电荷的场，不加实电荷自场。用 \(U=QV_{
\rm induced}/2\)
求力也可行，但求导时像位置与像大小都随真实电荷位置改变；若把它们冻结再乘1/2，会把力错误减半。独立能量求导与直接三项库仑力应使用同一
\(R\) 尺度，不能在中途把 \(a=\sqrt2R\) 的分母当成 \(R\)。

\subsection{六、关键洞察}\label{ux516dux5173ux952eux6d1eux5bdf}

两个边界的镜像操作必须成组出现。真实电荷与平面像电荷分别作球面镜像后，四个电荷关于
\(z=0\) 成反对称分布，因此平面仍保持零电势；而 Kelvin
镜像又分别保证球面零电势。

边界构造的检查顺序应是：先检查所有像电荷是否都在求解区域外，再逐个边界验证电势，最后才计算力。这个顺序能避免``公式算得很熟，但像电荷放进真实空间''的根本错误。

\textbf{镜像体系为何恰好闭合。} 设平面反射操作为 \(P\)，球面 Kelvin
反演操作为 \(K\)。本题真实电荷的位置满足到球心距离 \(2R\)，先作 \(P\)
再作 \(K\) 与先作 \(K\) 再作 \(P\)
都落在关于平面对称的两个内点，并且电量符号也相容。因此从真实电荷出发，经过``平面反射---球面反演''后只得到有限的四电荷体系。只要球心仍在接地平面上，平面反射与球面反演就可交换，对任意球外且平面上方的实电荷都闭合为这四电荷体系；闭合并不依赖本题距离恰好为
\(2R\)。改变两边界的相对位置或相交几何时才可能不再有限闭合。

\textbf{唯一性定理的作用。}
镜像法并不是在导体内部寻找真实电荷，而是在真实求解区域内构造一个满足泊松方程和全部边界电势的辅助场。只要像电荷均位于求解区域外，并且平面与半球面上都满足
\(V=0\)，静电唯一性定理就保证该辅助场等于真实场。因而``验证边界''是证明的一部分，不是计算后的可选检查。

\textbf{力与像电荷能量的区别。}
直接求力时，可以把真实电荷受到的导体作用等效为所有像电荷的库仑力之和；若改用静电能，则导体感应能量是

\[
U_{\rm ind}=\frac12Q\phi_{\rm image}(P),
\]

必须有因子
\(1/2\)。这是因为像电荷位置和电量随真实电荷位置一起变化，不能把真实电荷---像电荷系统当成两组相互独立的固定电荷。课堂上应把``力可直接相加''和``能量要乘一半''并列强调。

\textbf{平面边界的逐对验证。} 设真实电荷与平面像分别为

\[
Q:\ (a,0,b),
\qquad
-Q:\ (a,0,-b).
\]

对平面上任一点 \(X=(x,y,0)\)，两者到 \(X\) 的距离相等，所以

\[
\phi_{Q}(X)+\phi_{-Q}(X)=0.
\]

它们的 Kelvin 像分别位于

\[
\bm r_+=\frac{R^2}{a^2+b^2}(a,0,b),
\qquad
\bm r_- =\frac{R^2}{a^2+b^2}(a,0,-b),
\]

电量大小相等而符号相反。两位置仍关于 \(z=0\)
对称，因此第二对在平面上的电势也逐点抵消。这样不是只在交线或某几个特殊点满足零电势，而是在整个无限平面上满足边界条件。

\textbf{球面边界的逐对验证。} 对球外电荷 \(q\)，距球心 \(d>R\)，Kelvin
像为

\[
q'=-\frac Rdq,
\qquad
\bm r'=\frac{R^2}{d^2}\bm r_q.
\]

球面上任一点 \(X\) 满足几何恒等式

\[
|X-\bm r'|=\frac Rd|X-\bm r_q|.
\]

于是

\[
\frac{q}{|X-\bm r_q|}
+\frac{q'}{|X-\bm r'|}
=\frac{q}{|X-\bm r_q|}
-\frac{(R/d)q}{(R/d)|X-\bm r_q|}=0.
\]

真实电荷与自己的 Kelvin 像构成一组，平面像与它自己的 Kelvin
像构成另一组，两组分别使完整球面电势为零。半球只是该完整球面的一部分，因而自然满足要求。

\textbf{像电荷位置合法性的检查。}
本题真实求解区域在导体外部、平面上方。平面像位于 \(z<0\)；两个 Kelvin
像位于半径小于 \(R\)
的球内。三者均不在真实求解区域内，因此不会在外部区域引入虚假的泊松方程奇点。若某个像落到平面上方且球外，即使边界上数值暂时能抵消，也不能作为合法解。

\subsection{七、易错点}\label{ux4e03ux6613ux9519ux70b9}

\begin{itemize}
\tightlist
\item
  只放一个球面像电荷，破坏了导体平面的边界。
\item
  忘记平面像电荷的球面像电荷取正号。
\item
  把作用于真实电荷的力误算成真实电荷与全部四个电荷的合力。
\item
  在计算 \(q_3\) 的方向时把 \((3,0,5)\) 写成 \((5,0,3)\)。
\item
  把球面镜像的距离公式写成 \(r'=Rr_P\)，漏掉 \(R^2/r_P\) 的量纲。
\item
  计算力时把真实电荷本身产生的场也计入。
\item
  只验证半球弧面，不验证与它相接的无限平面。
\item
  用势能法时遗漏 \(U_{\rm ind}=Q\phi_{\rm image}/2\) 中的 \(1/2\)。
\item
  把像电荷当成导体内真实存在的离散电荷，进而错误讨论它们受到的力。
\end{itemize}

\subsection{八、课堂讲解建议}\label{ux516bux8bfeux5802ux8bb2ux89e3ux5efaux8bae}

先让学生分别说出``只有平面''和``只有完整球面''时的镜像，再要求检查每增加一个像电荷会不会破坏另一边界。最后用唯一性定理结束边界论证。

可要求学生先不算数值，只在图上标出每个像电荷的正负和大致位置，并预测
\(F_x,F_z\) 的符号。之后再做矢量代数，便能用预测检查计算。

\textbf{课堂推导顺序。} 第一步只处理平面，写出 \((a,0,-b)\) 处的
\(-Q\)；第二步把真实电荷和这个平面像电荷分别对完整球面作镜像；第三步让学生逐对检查
\(z=0\) 平面电势；第四步检查 \(r=R\)
球面电势；第五步才进入库仑力计算。每完成一个像电荷，都要回答``它解决了哪个边界、是否破坏另一个边界、它是否位于真实求解区域外''三个问题。

\textbf{板书计算组织。} 建议把三个像电荷按``电量---位置---到 \(P\)
的位移---距离---力分量''排成五列表格。先算无量纲分量，再统一乘
\(Q^2/(4\pi\varepsilon_0R^2)\)。这样可避免在每一行反复抄公共因子，也便于最后检查
\(F_y=0\)。

\textbf{理解检查。} 可问：若删去
\(+Q/2\)，哪一段边界首先失效？若把真实电荷移到 \(z\)
轴上，平面镜像仍需几个？为什么导体受到的总力不必与某一个像电荷力相等？这些问题能检验学生是否真正理解镜像法只是构造场的数学工具。

\textbf{一课时板书层次。}
第一栏只画几何和求解区域，明确真实电荷可活动的区域；第二栏写四个电荷的位置与符号；第三栏分别验证
\(z=0\) 与 \(r=R\)
上的电势；第四栏才列力分量。若学生基础较弱，可先只保留公共因子

\[
K=\frac{Q^2}{4\pi\varepsilon_0R^2},
\]

把所有坐标都除以
\(R\)，使距离计算变成无量纲向量运算。这样既减少根式干扰，也能让学生集中观察每个分量的符号。

\textbf{结果反查。} 得到合力后至少做四项检查：\(F_y=0\) 是否符合关于
\(xOz\) 面的对称性；量纲是否为 \(Q^2/(\varepsilon_0R^2)\)；靠近平面时
\(F_z\) 是否趋于吸引；把 \(z\) 坐标整体反号后，力的 \(z\)
分量是否相应反号。任何一项失败，都应先检查像电荷位置而非继续化简代数。

\textbf{分量符号的预判。} 平面像电荷 \(-Q\) 一定给出竖直向下的力；球内的
\(-Q/2\) 把真实电荷拉向球心附近，具有向左、向下分量；球内的 \(+Q/2\)
产生排斥，具有向右、向上分量。后两者距离不同，不能只按电量相等就判断抵消。正式计算前把三个箭头画在图
1-2 上，最终根式结果的每一项就有明确来源。

\textbf{教学时间分配。}
可用约三分之一课时复习单平面和单球面镜像，三分之一课时完成组合闭合与边界验证，最后三分之一课时计算力并做极限检查。若课堂时间有限，应保留边界验证而压缩根式化简，因为真正可迁移到新题的是镜像构造逻辑，而不是本题特定数字。

\subsection{九、变式训练}\label{ux4e5dux53d8ux5f0fux8badux7ec3}

\begin{enumerate}
\def\labelenumi{\arabic{enumi}.}
\tightlist
\item
  把真实电荷改放在 \(z\) 轴上，说明哪些力分量立即为零。
\item
  把半球改成半径为 \(R\) 的半球形凹坑，讨论镜像构造为什么不同。
\item
  将真实电荷位置写成
  \((a,0,b)\)，建立一般的四像电荷构造并说明何时能显著简化。
\end{enumerate}

\Needspace{3\baselineskip}
\textbf{变式解析要点}

\begin{enumerate}
\def\labelenumi{\arabic{enumi}.}
\tightlist
\item
  位于 \(z\) 轴上时，绕轴对称性先给出
  \(F_x=F_y=0\)；再判断该位置是否仍在半球外部，不能只凭对称性直接套本题四像体系。
\item
  凹坑问题的真实求解区域在球面另一侧，Kelvin
  像必须放在与原题不同的区域；判断像的位置比背电量公式更重要。
\item
  对一般 \((a,0,b)\)，先作平面像 \((a,0,-b)\)，再分别按
  \(R^2/\sqrt{a^2+b^2}\)
  的比例作球面像。本题平面通过球心，反射与反演可交换，对任意
  \(a^2+b^2>R^2\)、\(b>0\) 均形成四电荷闭合体系。
\end{enumerate}

\newpage

\section{第 2 题}\label{ux7b2c-2-ux9898}

半径为 \(r\)、质量为 \(m\) 的均匀薄壁圆筒在水平面上以速度 \(v\)
向右纯滚动。前方有质量同为 \(m\)、厚度为 \(h\)
的长垫块，垫块与地面间光滑，初始静止。圆筒与垫块尖角间摩擦足够大，碰撞完全非弹性；垫块长度足够大，始终不离开水平地面。

\begin{enumerate}
\def\labelenumi{\arabic{enumi}.}
\tightlist
\item
  当 \(h=r/4\) 时，求圆筒能不脱离垫块并翻上垫块的初速度范围。
\item
  求存在某个初速度使圆筒能不脱离垫块并翻上的 \(h\) 的范围。
\end{enumerate}

\begin{figure}[H]
\centering
\includegraphics[width=0.68\linewidth,height=\textheight,keepaspectratio,alt={图 2-1　纯滚动圆筒与可水平移动垫块的碰撞几何。}]{../figures/question-02.pdf}
\caption*{图 2-1　纯滚动圆筒与可水平移动垫块的碰撞几何。}
\end{figure}

\subsection{一、题型定位}\label{ux4e00ux9898ux578bux5b9aux4f4d-1}

碰撞先决定垫块平动与圆筒转动，越障再检查能量和接触力。垫块一直在加速，法向方程必须包含它的运动。

\subsection{二、知识点}\label{ux4e8cux77e5ux8bc6ux70b9-1}

\begin{itemize}
\tightlist
\item
  薄壁圆筒转动惯量 \(I_C=mr^2\)。
\item
  完全非弹性碰撞的共同接触点速度。
\item
  瞬时冲量下的角动量守恒。
\item
  支持力临界条件。
\item
  机械能与角动量在越障阶段的配合。
\end{itemize}

\subsection{三、物理图像}\label{ux4e09ux7269ux7406ux56feux50cf-1}

碰撞瞬间圆筒与尖角黏住，接触点 \(P\)
成为圆筒相对垫块的瞬时转动中心，但垫块本身还会平动。碰撞后若支持力为负，圆筒立即脱离；若速度过小，则虽不脱离也翻不到垫块上表面。

\subsection{四、解题入口}\label{ux56dbux89e3ux9898ux5165ux53e3-1}

令碰后垫块速度为 \(u\)，圆筒相对垫块绕 \(P\) 的角速度为
\(\omega\)。先由几何关系确定 \(\theta\)，再列水平动量和关于 \(P\)
的角动量方程解出 \(u,\omega\)。

\subsection{五、解析}\label{ux4e94ux89e3ux6790-1}

\begin{figure}[H]
\centering
\includegraphics[width=0.9\linewidth,height=\textheight,keepaspectratio,alt={图 2-2　碰后以尖角为瞬时转轴的几何与速度变量。}]{../figures/solution-02.pdf}
\caption*{图 2-2　碰后以尖角为瞬时转轴的几何与速度变量。}
\end{figure}

图 2-2 将圆心速度拆为随垫块平动的 \(u\hat{\bm x}\) 与绕尖角 \(P\) 转动的
\(\bm\omega\times\bm r_{PC}\)，并标出高度关系中的角
\(\theta\)。下面的水平动量、关于 \(P\)
的角动量以及支持力条件都按这组正方向列式。

设圆心到尖角的连线与水平线夹角为 \(\theta\)。由竖直高度差

\[r\sin\theta=r-h,\]

故

\[\boxed{\sin\theta=1-\frac hr}.\]

薄壁圆筒初始角速度为 \(v/r\)。

\Needspace{3\baselineskip}
\textbf{1. 碰撞后的速度}

系统不受水平外冲量，故水平动量守恒。圆筒质心速度由垫块平动 \(u\) 与绕
\(P\) 的转动叠加，其水平分量为 \(u+r\omega\sin\theta\)。于是

\[
2mu+mr\omega\sin\theta=mv. \tag{1}
\]

碰撞冲量作用于 \(P\)，其对 \(P\)
的力矩为零。只对圆筒取碰撞瞬间固定空间位置 \(P\)
为矩心；尖角对圆筒的冲量矩为零。垫块的地面竖直冲量对 \(P\)
可能有矩，不能声称整个系统关于 \(P\)
的角动量都守恒。碰前圆筒质心平动贡献为 \(mvr\sin\theta\)，自转贡献为
\(mr^2(v/r)=mrv\)。碰后角动量包括圆筒质心随 \(u\)
的平动、圆筒绕质心自转以及质心绕 \(P\) 的转动：

\[
mvr\sin\theta+mrv
=mur\sin\theta+mr^2\omega+mr^2\omega. \tag{2}
\]

约去 \(m\)，联立 (1)(2)，得到

\[
\boxed{u=\frac{2-\sin\theta-\sin^2\theta}{4-\sin^2\theta}v},
\]

\[
\boxed{\omega=\frac{2+\sin\theta}{4-\sin^2\theta}\frac vr}.
\]

\Needspace{3\baselineskip}
\textbf{2. 碰后完整动力学}

取垫块水平坐标 \(X\)、角度 \(\theta\)
从圆心到尖角连线的水平投影计起，圆心坐标为

\[x_C=X-r\cos\theta,\quad y_C=h+r\sin\theta.\]

无滑动时圆筒自转速率为 \(\dot\theta\)，记
\(w=\dot\theta>0\)。总动能与水平动量是

\[K=m\dot X^2+mr\sin\theta\,\dot X w+mr^2w^2,\]

\[p_x=2m\dot X+mr\sin\theta w=mv.\]

消去 \(\dot X\) 后

\[K=\frac{p_x^2}{4m}+\frac{mr^2}{4}(4-\sin^2\theta)w^2.\]

由能量与角方程，设 \(s=\sin\theta\)、\(D=4-s^2\)，有

\[Dw^2=D_0w_0^2-\frac{4g}{r}(s-s_0),\]

\[\ddot\theta=\frac{\cos\theta(s w^2-2g/r)}{D},\quad
\ddot X=-\frac r2(\cos\theta w^2+s\ddot\theta).\]

式中的基底加速度不能删去。沿从尖角指向圆心的单位向量投影，法向接触力为

\[\boxed{\frac Nm=\frac{3g s-r(2+s^2)w^2}{4-s^2}}.\tag{3}\]

原稿 \(N=m(gs-rw^2)\) 把移动垫块当作惯性参考系，漏掉
\(\ddot X\)，因此需修正。

\Needspace{3\baselineskip}
\textbf{3. 能量下限与全过程接触上限}

顶部 \(s=1\)，只令相对转动趋零，系统仍保留
\(p_x^2/(4m)\)。越障能量下界为

\[w_{0,\inf}^2=\frac{4gh}{r^2(4-s_0^2)}.\]

初始法向力非负给出

\[w_0^2\le\frac{3g s_0}{r(2+s_0^2)}.\]

这不只是瞬间条件。消去 \(w\) 后，全程 \(N\ge0\) 等价于

\[\frac r gD_0w_0^2+4s_0\le f(s),\qquad
f(s)=\frac{s(20+s^2)}{2+s^2}.\]

在 \(s_0\le s\le1\) 内

\[f'(s)=\frac{40-14s^2+s^4}{(2+s^2)^2}>0,\]

故最严格处确为刚碰后。切向摩擦力为
\(T=-mr\ddot\theta\)，对任何固定有限摩擦因数仍须满足
\(|T|\le\mu N\)；原题``摩擦足够大''按理想无滑动处理，但上界 \(N=0\)
时一般 \(T\ne0\)，不能由有限摩擦精确取到。

\Needspace{3\baselineskip}
\textbf{4. \(h=r/4\) 的修正区间}

\(s_0=3/4\)、\(w_0=4v/(5r)\)，能量下限保留，接触上限改为

\[\boxed{\sqrt{\frac{5gr}{11}}<v<\sqrt{\frac{225gr}{164}}}.\]

这是理想无滑动、实际在有限时间翻过顶部的开区间。下界只渐近达到顶部，上界要求零法向力下仍有切向力；若只列形式临界包络，可写闭区间并明确两端不能作为有限摩擦下的实际翻越。

\Needspace{3\baselineskip}
\textbf{5. 厚度范围}

令 \(s_0=1-h/r\)，上下界有严格重叠要求

\[4(1-s_0)(2+s_0^2)<3s_0(4-s_0^2),\]

即

\[s_0^3-4s_0^2+20s_0-8>0.\]

该多项式在 \([0,1]\) 严格递增，唯一根
\(s_c=0.4335131439\)。因此理想无滑动的厚度范围为

\[\boxed{0<h<0.5664868561r}.\]

原稿 \(0.5665r\) 来自漏掉基底加速度的法向条件。有限指定 \(\mu\)
会进一步缩小可行域；题目未给 \(\mu\)，这里只能给足够大摩擦的理想包络。

\Needspace{3\baselineskip}
\textbf{碰撞阶段为什么不能守恒机械能}

圆筒与尖角发生完全非弹性碰撞，接触处存在不可逆形变和摩擦冲量。碰撞前后动能差转化为内能，故

\[K_{\rm after}<K_{\rm before}.\]

但地面对系统的外冲量竖直，水平动量仍守恒；接触冲量通过 \(P\)，对 \(P\)
的冲量矩为零，所以圆筒自身对碰撞瞬间 \(P\)
的角动量守恒。这正是两个守恒式可以使用而机械能式不能使用的原因。

\Needspace{3\baselineskip}
\textbf{越障下限的完整消元}

碰后圆筒相对尖角的转动惯量包括绕质心和质心绕 \(P\) 两部分：

\[I_P=I_C+mr^2=2mr^2.\]

临界翻上时圆筒相对垫块的角速度降为零，但系统因水平动量仍以共同速度
\(u_2\) 平动。碰后初态到临界末态的能量关系写成

\[
\frac12I_P\omega^2
=\frac12(2m)u_2^2+mgh. \tag{5}
\]

同一过程中水平动量给出

\[mr\omega\sin\theta=2mu_2,\]

故

\[u_2=\frac12r\omega\sin\theta.\]

代入 (5)：

\[
mr^2\omega^2
=m\frac{r^2\omega^2\sin^2\theta}{4}+mgh.
\]

约去 \(m\)，移项：

\[
r^2\omega^2\left(1-\frac{\sin^2\theta}{4}\right)=gh.
\]

所以

\[
\omega_2^2=\frac{4gh}{r^2(4-\sin^2\theta)},
\]

这正是正文使用的式 (4)。

\Needspace{3\baselineskip}
\textbf{两个临界速度的物理含义}

\(\omega<\omega_2\) 时，相对转动动能不足以克服
\(mgh\)；\(\omega>\omega_1\)
时，圆筒所需向心力超过重力在径向所能提供的分量，接触力必须变成拉力才可维持，实际就会脱离。因此成功翻越只可能出现在两临界值之间。

\Needspace{3\baselineskip}
\textbf{碰撞方程的矩阵解法}

方程 (1)(2) 可写成

\[
\begin{pmatrix}
2&r\sin\theta\\
r\sin\theta&2r^2
\end{pmatrix}
\begin{pmatrix}u\\\omega\end{pmatrix}
=
\begin{pmatrix}v\\rv(1+\sin\theta)\end{pmatrix}.
\]

系数行列式

\[D=r^2(4-\sin^2\theta)>0,\]

保证物理几何范围内碰后速度唯一。克拉默法则逐项展开就得到正文中的 \(u\)
和 \(\omega\)，也说明分母 \(4-\sin^2\theta\)
来自耦合平动与转动的线性方程组。

\Needspace{3\baselineskip}
\textbf{临界厚度附近的区间收缩}

当 \(h\) 增大到 \(0.5665r\)
时，\(\omega_1=\omega_2\)，允许区间收缩成一个临界速度。再增厚，翻越所需的能量下限超过保持接触所允许的速度上限，因而不存在任何成功初速度。

\Needspace{3\baselineskip}
\textbf{碰撞能量的独立验算与摩擦边界}

由两条碰撞方程可化得

\[K_+/K_-=(2+s_0)/(4-s_0^2),\qquad K_-=mv^2.\]

在 \(0\le s_0<1\) 时该比值小于 \(1\)，\(s_0\to1\) 无台阶极限趋于
\(1\)。必须包含圆心竖直动能 \(m(rw\cos\theta)^2/2\)，才会得到此式。

由 \(T=-mr\ddot\theta\) 可算每一时刻所需
\(\mu_{\rm req}=|T|/N\)。在上述开区间且碰后至顶部 \(N\)
始终正的情况下，这个比值有有限最大值，因而存在``足够大''的有限摩擦；向上边界逼近时最大比值发散。把支持力、能量和摩擦三项分开核对，能区分形式临界包络与真实可实现轨迹。

顶部临界不能有限时间达到：在恰等能量时，\(1-\sin\theta\sim(\pi/2-\theta)^2/2\)，故
\(w\propto\pi/2-\theta\)，积分 \(dt=d\theta/w\)
发散。严格翻越必须超过能量下界。

\Needspace{3\baselineskip}
\textbf{碰撞的两条方程怎样解出 \(u,w_0\)}

令 \(s_0=1-h/r\)、\(c_0=\sqrt{1-s_0^2}\)，垫块碰后水平速度为
\(u\)。无滑动接触给圆心速度

\[\bm v_C=(u+rs_0w_0,rc_0w_0),\]

圆筒自转为顺时针 \(w_0\)。总水平动量守恒为

\[2u+rs_0w_0=v.\]

只对圆筒取瞬时 \(P\) 为矩心，碰前角动量绝对值为 \(mrv(1+s_0)\)，碰后为
\(mr(s_0u+2rw_0)\)。因此

\[s_0u+2rw_0=(1+s_0)v.\]

将第一式乘 \(s_0/2\) 后从第二式减去，得到

\[w_0=\frac{2+s_0}{4-s_0^2}\frac vr,\qquad
u=\frac{2-s_0-s_0^2}{4-s_0^2}v.\]

在 \(s_0=3/4\) 时为 \(w_0=0.8v/r,u=0.2v\)，圆心水平速度
\(u+rs_0w_0=0.8v\)，竖直速度
\(rc_0w_0>0\)。解析图中圆心箭头向右上、垫块箭头水平向右，分别对应这两个速度，不能画成碰后两质心同速。

\subsection{六、关键洞察}\label{ux516dux5173ux952eux6d1eux5bdf-1}

碰撞和越障是两个阶段，守恒量不同。碰撞阶段机械能不守恒，却可用水平动量和对冲量作用点的角动量；越障阶段则应从碰后状态重新讨论支持力、角动量和能量。

速度区间的上下界来自完全不同的失效机制：能量不足给下界，保持接触给上界。看到``既不能太慢也不能太快''时，应主动寻找两个独立临界条件。

\textbf{接触点约束的准确含义。}
``完全非弹性''只说明碰撞瞬间圆筒与尖角的相对速度在接触法向与切向上满足黏着约束，并不表示碰后圆筒与垫块永远焊在一起。碰后若维持接触所需的支持力变成负值，单侧接触会自动解除。因此碰撞方程给出的只是初始状态，随后仍必须单独检查
\(N\ge0\)。

\textbf{为什么水平动量始终可用。}
地面对垫块的支持冲量以及重力冲量均竖直，系统在水平方向没有外冲量；碰撞后地面仍光滑，也没有水平外力。因此水平动量既能用于碰撞瞬间，也能用于后续越障过程。但对尖角
\(P\) 的角动量守恒只适用于碰撞极短时间，碰后 \(P\) 随垫块运动，且重力对
\(P\) 有力矩，不能把同一角动量守恒式延用到整个翻越过程。

\textbf{能量下限的结构。}
越障阶段机械能守恒并不意味着``碰后全部动能都能升高圆筒''。系统水平总动量非零，因此系统质心的平动动能不能消失。最稳妥的做法是把动能分为质心平动与相对运动，只有相对运动部分能够在临界态转化为
\(mgh\)。这正是式 (4) 分母出现 \(4-\sin^2\theta\) 的物理来源。

\Needspace{3\baselineskip}
\textbf{支持力上限的完整列式。}

图 2-2 区分圆心速度 \(\bm v_C\) 和尖角速度
\(\dot X\hat x\)，并在圆筒上画出法向 \(N\)、切向 \(T\) 和
\(mg\)。圆心径向加速度为 \(-r w^2-\cos\theta\ddot X\)，故

\[N-mg\sin\theta=m[-r w^2-\cos\theta\ddot X].\]

代入水平动量的导数与角运动方程，化为式（3）。这一过程同时包含平动和转动，不允许把随垫块坐标系误当惯性系。

\textbf{碰撞方程逐项解释。} 水平动量式

\[
2mu+mr\omega\sin\theta=mv
\]

中的 \(2mu\) 来自垫块与圆筒共同随基底平动，\(mr\omega\sin\theta\)
是圆筒相对尖角转动速度的水平分量。关于 \(P\)
的角动量式中，碰前有质心平动 \(mvr\sin\theta\) 与自转
\(mrv\)；碰后有随垫块平动的 \(mur\sin\theta\)、绕质心自转 \(mr^2\omega\)
和质心绕 \(P\) 的 \(mr^2\omega\)。要求学生为每一项在图 2-2
上找到对应箭头或力臂，可避免凭印象凑方程。

\textbf{碰撞能量损失的定性核验。} 由联立方程得到 \(u\) 与 \(\omega\)
后，可以把碰后动能写成

\[
K_+=\frac12mu^2
+\frac12m\left(u+r\omega\sin\theta\right)^2
+\frac12m(r\omega\cos\theta)^2+\frac12mr^2\omega^2.
\]

碰前动能为薄壁圆筒的平动与转动之和

\[
K_-=\frac12mv^2+\frac12mr^2\left(\frac vr\right)^2=mv^2.
\]

代入任意 \(0\le\sin\theta\le1\) 的数值都应得到
\(K_+<K_-\)。此核验不用于求未知量，但能检查角速度符号和动能漏项；若出现
\(K_+>K_-\)，碰撞方程一定有误。

\subsection{七、易错点}\label{ux4e03ux6613ux9519ux70b9-1}

\begin{itemize}
\tightlist
\item
  把垫块当作固定台阶，漏掉速度 \(u\)。
\item
  在完全非弹性碰撞中使用机械能守恒。
\item
  只检查翻越能量，不检查刚碰后的支持力。
\item
  把 \(h/r\) 直接当成 \(\sin\theta\)，正确关系是 \(\sin\theta=1-h/r\)。
\item
  碰后把 \(P\) 当作地面固定点，忘记它随垫块以速度 \(u\) 移动。
\item
  在越障末态令整个系统静止，违反水平动量守恒。
\item
  把``完全非弹性''误解成碰后永久黏住，从而不再检查脱离条件。
\item
  认为关于 \(P\) 的角动量在整个越障过程中都守恒。
\item
  只比较碰前和最高点能量，忽略碰撞阶段已经损失的机械能。
\end{itemize}

\subsection{八、课堂讲解建议}\label{ux516bux8bfeux5802ux8bb2ux89e3ux5efaux8bae-1}

把板书分成``碰撞前---碰撞瞬间---碰撞后越障''三栏，每一栏只写这一阶段允许使用的守恒定律。用速度区间的上下限分别解释``太慢''和``太快''的失败机制。

推导 \(\omega_2\) 时先保留末态共同平动速度
\(u_2\)，用水平动量消去它。这个步骤正好展示为什么移动台阶与固定台阶的临界速度不同。

\textbf{课堂推导顺序。} 先完成几何关系
\(\sin\theta=(r-h)/r\)，再把碰撞阶段框起来，只列水平动量和关于 \(P\)
的角动量；求得 \(u,\omega\)
后画一条时间分界线，明确后续讨论已进入碰后连续运动。接着分别用支持力临界和越障能量临界得到上下限，最后才代入
\(h=r/4\)。

\textbf{为什么选择 \(P\) 点。}
若关于圆心列角动量，未知冲量会产生力矩；若关于地面某固定点列角动量，垫块和圆筒的表达式都更复杂。选择冲量作用点
\(P\)，正是为了消去未知碰撞冲量。这是竞赛碰撞题中``选矩心''的通用策略。

\textbf{诊断学生错误。}
若学生得到单一临界速度，多半只做了能量或支持力其中之一；若结果与垫块质量无关，多半把垫块当固定；若碰后速度比碰前更大且又声称机械能守恒，应要求其直接计算碰前后动能差并查找阶段混淆。

\textbf{速度图的使用。} 在解析图上把圆心速度拆成两部分：随垫块的
\(u\hat{\bm x}\) 与绕 \(P\) 的
\(\bm\omega\times\bm r_{PC}\)。要求学生先用图判断水平分量的正负，再写
\(v_{Cx}=u+r\omega\sin\theta\)。这样可以避免把 \(\sin\theta\) 与
\(\cos\theta\)
互换。碰后角速度的正方向也应在图上一次性规定，后续各式沿用同一约定。

\textbf{临界状态的三幅图。}
建议依次画``刚碰后''\,``刚好不脱离''\,``刚好到达上表面''三幅小图。第一幅标
\(u,\omega\)；第二幅标 \(N=0\)
和径向重力分量；第三幅标相对角速度为零、但系统仍有共同平动速度
\(u_2\)。三幅图能直观纠正``最高点全系统静止''的常见误解。

\textbf{结果检查。} 当 \(h\to0\) 时，势垒消失，能量下限应趋于零；当
\(h\)
增大时，下限上升、接触上限下降，成功区间应收缩；当垫块质量趋于无穷时，\(u\to0\)，应回到固定尖角模型。这三项比单纯代数验算更能暴露建模错误。

\textbf{厚度临界的图像讲法。} 在同一坐标图上画 \(\omega_1(h)\) 与
\(\omega_2(h)\)：前者表示接触允许的上限，后者表示越障要求的下限。小厚度时两曲线之间有可行带；随
\(h\) 增大，可行带逐渐缩窄；交点就是
\(h_c\approx0.5665r\)。这比直接展示三次方程的根更能说明``临界厚度为何存在''。

\textbf{三次不等式的筛选。} 令 \(x=1-h/r\) 后，物理范围为
\(0\le x<1\)。虽然 \(x^3-4x^2+20x-8=0\)
在实数域可能有多个根，但只有落在该物理区间的根才决定厚度。先用
\(f'(x)=3x^2-8x+20>0\)
证明区间内严格单调，再用端点异号确定唯一根，避免计算器列出多个根后凭大小猜选。

\textbf{课堂练习的分层。}
基础层只要求写出几何关系和两条碰撞方程；进阶层推导 \(N\ge0\)
与越障能量；竞赛层完成厚度临界三次不等式。三层共用同一张图和同一套变量，便于教师根据学生状态在课堂中调整深度，而不改变主线。

\subsection{九、变式训练}\label{ux4e5dux53d8ux5f0fux8badux7ec3-1}

\begin{enumerate}
\def\labelenumi{\arabic{enumi}.}
\tightlist
\item
  若垫块质量改为 \(M\)，重新建立两条碰撞方程。
\item
  若圆筒为实心圆柱，指出哪些系数因 \(I_C=mr^2/2\) 而改变。
\item
  令垫块质量趋于无穷，验证结果能否回到固定台阶模型。
\end{enumerate}

\Needspace{3\baselineskip}
\textbf{变式解析要点}

\begin{enumerate}
\def\labelenumi{\arabic{enumi}.}
\tightlist
\item
  质量改为 \(M\) 后，水平动量式中的 \(2mu\) 改为 \((M+m)u\)；关于瞬间
  \(P\)
  的角冲量式只对圆筒写，仍不含垫块项；垫块质量通过总水平动量方程进入。
\item
  实心圆柱只改变自转惯量，碰前自转角动量和碰后绕质心部分都应把 \(mr^2\)
  换成 \(mr^2/2\)；质心绕 \(P\) 的 \(mr^2\) 不变。
\item
  令 \(M/m\to\infty\) 后先检查
  \(u\to0\)，再把剩余角速度与固定台阶碰撞结果比较。若只在最终速度公式中令
  \(M\to\infty\) 而没有检查守恒式，容易掩盖漏项。
\end{enumerate}

\newpage

\section{第 3 题}\label{ux7b2c-3-ux9898}

一根均匀弹簧置于无摩擦水平管中，管子绕一端的竖直轴以恒定角速度
\(\omega\) 转动。弹簧内端固定在距转轴 \(r_0\) 处。弹簧质量为
\(m\)、劲度系数为 \(k\)、自然长度为
\(l_0\)，遵从胡克定律且不会损坏。求稳定伸长后的长度
\(l\)，说明解存在的约束；并讨论 \(r_0\to0\) 的情况。

\begin{figure}[H]
\centering
\includegraphics[width=0.82\linewidth,height=\textheight,keepaspectratio,alt={图 3-1　随径向管匀速转动的有质量弹簧。}]{../figures/question-03.pdf}
\caption*{图 3-1　随径向管匀速转动的有质量弹簧。}
\end{figure}

\subsection{一、题型定位}\label{ux4e00ux9898ux578bux5b9aux4f4d-2}

用自然材料坐标标记质量，用变形后的半径写离心力。张力由自由端向内累积，解出形函数后要检查自由端零张力和低速极限。

\subsection{二、知识点}\label{ux4e8cux77e5ux8bc6ux70b9-2}

\begin{itemize}
\tightlist
\item
  均匀弹簧的局部胡克定律。
\item
  线密度与材料坐标。
\item
  旋转参考系中的离心力。
\item
  二阶常系数微分方程。
\item
  边界条件与临界角速度。
\end{itemize}

\subsection{三、物理图像}\label{ux4e09ux7269ux7406ux56feux50cf-2}

越靠近外端的微元离转轴越远，需要的向心力越大；内侧截面还要拉住它外面的全部质量，因此弹簧内外各处伸长率不同。自由端张力为零，固定端位置给定。

\subsection{四、解题入口}\label{ux56dbux89e3ux9898ux5165ux53e3-2}

用自然状态的材料坐标 \(x\in[0,l_0]\) 标记微元，令 \(r(x)\)
为该微元转动后到转轴的距离。局部伸长率为 \(dr/dx\)，张力由它决定。

\subsection{五、解析}\label{ux4e94ux89e3ux6790-2}

\begin{figure}[H]
\centering
\includegraphics[width=0.9\linewidth,height=\textheight,keepaspectratio,alt={图 3-2　弹簧微元的张力与离心力平衡。}]{../figures/solution-03.pdf}
\caption*{图 3-2　弹簧微元的张力与离心力平衡。}
\end{figure}

图 3-2 用自然材料坐标 \(x\) 标记微元 \(dx\)，而 \(r(x)\)
表示变形后到转轴的空间位置；两者不能混用。以下张力差 \(F(x+dx)-F(x)\)
和离心力 \(dm\,\omega^2r\) 将按图示方向逐项平衡。

自然状态线密度为

\[\mu_0=\frac{m}{l_0}.\]

材料微元 \(dx\) 的质量为

\[dm=\mu_0dx=\frac{m}{l_0}dx.\]

\Needspace{3\baselineskip}
\textbf{1. 局部胡克定律}

若整个弹簧均匀伸长 \(\Delta l\)，张力是 \(k\Delta l\)。对自然长度 \(dx\)
的微段，其实际长度为 \(dr\)，伸长为 \(dr-dx\)。该微段等效劲度系数为
\(kl_0/dx\)，因此

\[
F(x)=\frac{kl_0}{dx}(dr-dx)
=kl_0\left(\frac{dr}{dx}-1\right). \tag{1}
\]

\Needspace{3\baselineskip}
\textbf{2. 微元平衡}

在随管转动的参考系中，微元受到离心力

\[dm\,\omega^2r=\frac{m}{l_0}\omega^2r,dx.\]

两端张力之差满足

\[F(x+dx)-F(x)=-\frac{m}{l_0}\omega^2r,dx.\]

除以 \(dx\) 并令 \(dx\to0\)：

\[\frac{dF}{dx}=-\frac{m}{l_0}\omega^2r.\]

代入 (1)：

\[kl_0\frac{d^2r}{dx^2}=-\frac{m}{l_0}\omega^2r.\]

定义

\[
\beta=\omega\sqrt{\frac{m}{k}},
\]

并改用无量纲材料坐标 \(\xi=x/l_0\)，则

\[\frac{d^2r}{d\xi^2}+\beta^2r=0. \tag{2}
\]

\Needspace{3\baselineskip}
\textbf{3. 边界条件}

固定端位置给定：

\[r(0)=r_0. \tag{3}
\]

自由端张力为零。由 (1)，有

\[F(l_0)=0
\quad\Rightarrow\quad
\left.\frac{dr}{dx}\right|_{x=l_0}=1.
\]

换成 \(\xi\)：

\[\left.\frac{dr}{d\xi}\right|_{\xi=1}=l_0. \tag{4}
\]

方程 (2) 的通解为

\[r(\xi)=C\cos(\beta\xi)+D\sin(\beta\xi).\]

由 (3) 得 \(C=r_0\)。对它求导：

\[
\frac{dr}{d\xi}
=-r_0\beta\sin(\beta\xi)+D\beta\cos(\beta\xi).
\]

代入 (4)：

\[-r_0\beta\sin\beta+D\beta\cos\beta=l_0,\]

所以

\[
D=\frac{l_0/\beta+r_0\sin\beta}{\cos\beta}.
\]

\Needspace{3\baselineskip}
\textbf{4. 总长度}

外端位置为 \(r(1)\)，故弹簧长度

\[l=r(1)-r_0.\]

逐步代入：

\[
\begin{aligned}
r(1)
&=r_0\cos\beta
+\frac{l_0/\beta+r_0\sin\beta}{\cos\beta}\sin\beta\\
&=\frac{l_0}{\beta}\tan\beta
+r_0\left(\cos\beta+\frac{\sin^2\beta}{\cos\beta}\right)\\
&=\frac{l_0}{\beta}\tan\beta+r_0\sec\beta.
\end{aligned}
\]

因此

\[
\boxed{
l=\frac{l_0}{\beta}\tan\beta+r_0(\sec\beta-1)
},
\qquad
\beta=\omega\sqrt{\frac mk}.
\]

\Needspace{3\baselineskip}
\textbf{5. 存在条件与 \(r_0\to0\)}

要使 \(r(\xi)\) 保持有限、伸长率为正且不跨过转轴，第一物理解支要求

\[0\le\beta<\frac\pi2.\]

即

\[
\boxed{\omega<\frac\pi2\sqrt{\frac km}}.
\]

当 \(r_0\to0\) 时

\[
\boxed{
l=\frac{l_0}{\beta}\tan\beta
=l_0\sqrt{\frac{k}{m}}\frac1\omega
\tan\left(\omega\sqrt{\frac mk}\right)
}.
\]

\Needspace{3\baselineskip}
\textbf{低速展开核验}

利用
\(\tan\beta/\beta=1+\beta^2/3+O(\beta^4)\)、\(\sec\beta-1=\beta^2/2+O(\beta^4)\)，得到

\[
l=l_0+\omega^2\frac mk
\left(\frac{l_0}{3}+\frac{r_0}{2}\right)+O(\omega^4).
\]

当 \(\omega\to0\) 时 \(l\to l_0\)，符合静止极限。

\Needspace{3\baselineskip}
\textbf{张力分布的显式表达}

由 \(F=kl_0(dr/dx-1)\) 以及 \(\xi=x/l_0\)，有

\[
F(\xi)=k\left(\frac{dr}{d\xi}-l_0\right).
\]

代入 \(r(\xi)=r_0\cos(\beta\xi)+D\sin(\beta\xi)\)，得到

\[
F(\xi)=k\left[-r_0\beta\sin(\beta\xi)
+D\beta\cos(\beta\xi)-l_0\right].
\]

在 \(\xi=1\) 处由边界条件确有
\(F(1)=0\)；向内移动时张力逐渐增大，符合内侧截面需要拉住更多外部质量的物理图像。

\Needspace{3\baselineskip}
\textbf{离散模型到连续方程}

可把弹簧分成 \(n\) 段，每段质量 \(m/n\)、自然长度 \(l_0/n\)、劲度系数
\(nk\)。第 \(j\) 个截面张力之差为

\[F_{j+1}-F_j=-\frac mn\omega^2r_j.\]

又有

\[F_j=nk\left[(r_j-r_{j-1})-\frac{l_0}{n}\right].\]

联立后

\[
nk(r_{j+1}-2r_j+r_{j-1})
=-\frac mn\omega^2r_j.
\]

两边乘 \(n/l_0^2\) 并令 \(n\to\infty\)，差分二阶导数趋于
\(d^2r/dx^2\)，就得到正文的连续方程。这个过程说明微分方程并非凭空写出。

\Needspace{3\baselineskip}
\textbf{临界角速度的物理解释}

当 \(\beta\to\pi/2^-\) 时，\(\cos\beta\to0^+\)，自由端边界要求系数 \(D\)
发散，因而长度与固定端张力同时发散。理想胡克弹簧模型无法在更高角速度维持有限稳态；真实弹簧会在接近该点前进入非线性或损坏。

\Needspace{3\baselineskip}
\textbf{形函数的边界核验}

完整形函数是

\[
r(\xi)=r_0\cos(\beta\xi)
+\frac{l_0/\beta+r_0\sin\beta}{\cos\beta}
\sin(\beta\xi).
\]

令 \(\xi=0\) 得 \(r(0)=r_0\)。求导并令 \(\xi=1\)：

\[
r'(1)=-r_0\beta\sin\beta
+\left(l_0/\beta+r_0\sin\beta\right)\beta=l_0.
\]

所以 \(F(1)=k[r'(1)-l_0]=0\)，两个边界条件均逐项满足。

\Needspace{3\baselineskip}
\textbf{固定端张力的低速极限}

由 \(F(0)=k[r'(0)-l_0]\) 得

\[
F(0)=k\left[
\frac{l_0+r_0\beta\sin\beta}{\cos\beta}-l_0
\right].
\]

低速展开为

\[F(0)\simeq m\omega^2\left(\frac{l_0}{2}+r_0\right),\]

这个结果还可直接用低速极限下的整体受力复核。此时弹簧伸长是二阶小量，可在计算离心力时先取
\(r(x)\simeq r_0+x\)。线密度为 \(m/l_0\)，位于 \(x\) 到 \(x+dx\)
的微元所需向心力大小为

\[
dF_c=\frac{m}{l_0}\omega^2(r_0+x)\,dx.
\]

固定端必须拉住其外侧的全部微元，故

\[
\begin{aligned}
F(0)
&=\int_0^{l_0}dF_c\\
&=\frac{m\omega^2}{l_0}
\int_0^{l_0}(r_0+x)\,dx\\
&=\frac{m\omega^2}{l_0}
\left(r_0l_0+\frac{l_0^2}{2}\right)\\
&=m\omega^2\left(r_0+\frac{l_0}{2}\right).
\end{aligned}
\]

它与精确解的低速展开完全相同。这个核验还说明两项的来源：\(m\omega^2r_0\)
是整根弹簧随固定端半径平移所需的部分，\(m\omega^2l_0/2\)
是弹簧自身沿径向分布所需的部分。

\Needspace{3\baselineskip}
\textbf{形函数的稳定性检验}

对任意满足固定端 \(\eta(0)=0\) 的小扰动，旋转系有效势的二次变分为

\[\delta^2 U=kl_0\int_0^{l_0}(\eta_x)^2dx
-\frac{m\omega^2}{l_0}\int_0^{l_0}\eta^2dx.\]

固定---自由端的最小本征值为 \((\pi/2l_0)^2\)，所以二次型正定恰需
\(\omega^2<\pi^2 k/(4m)\)，与首个正切奇点一致。这里稳定表示小径向扰动不会指数增长；无阻尼管道不会让任意初态自动停在这个形状。

图 3-2 左端的张力 \(F(x)\) 向内，右端张力 \(F(x+dx)\)
向外，离心力向外，三者和为零。原图把离心力画向内、且用变形后长度作质量坐标，均已纠正。

\Needspace{3\baselineskip}
\textbf{质量坐标能避免把伸长后的密度当常量}

取弹簧原长坐标 \(0\le x\le l_0\)，线密度始终是 \(m/l_0\)。平衡方程

\[kl_0r_{xx}=-\frac m{l_0}\omega^2r\]

化为 \(r_{xx}+(\beta/l_0)^2r=0\)，\(\beta=\omega\sqrt{m/k}\)。固定端
\(r(0)=r_0\)，自由端张力为零给 \(r_x(l_0)=1\)。通解

\[r(x)=r_0\cos(\beta x/l_0)+A\sin(\beta x/l_0),\]
\[A=\frac{l_0/\beta+r_0\sin\beta}{\cos\beta}.\]

把 \(x=l_0\) 代入并减去 \(r_0\)，得到
\(l=l_0\tan\beta/\beta+r_0(\sec\beta-1)\)。小转速展开

\[l=l_0+\beta^2(l_0/3+r_0/2)+O(\beta^4),\]

伸长为正、随 \(\omega^2\) 起步，检查了符号与静态极限。用 \(m/l\)
作为材料坐标密度会漏掉质量守恒的雅可比，图中已改为 \(dm=(m/l_0)dx\)。

\Needspace{3\baselineskip}
\textbf{端点加质量后的边界与首个临界值}

在自由端加质量 \(M\)，设
\(\eta=M/m\)、\(\beta=\omega\sqrt{m/k}\)。微元方程不变，仍可写

\[r(x)=r_0\cos(\beta x/l_0)+B\sin(\beta x/l_0).\]

端点质量要求张力提供向心力，边界改为

\[kl_0[r_x(l_0)-1]=M\omega^2r(l_0).\]

代入形函数，得到

\[B=\frac{l_0+\beta r_0\sin\beta+\eta\beta^2r_0\cos\beta}
{\beta\cos\beta-\eta\beta^2\sin\beta}.\]

实际弹簧长度为 \(l=r(l_0)-r_0\)。令
\(M\to0\)，该式还原原自由端解。首个分母零点满足

\[\cot\beta_c=\eta\beta_c.\]

对 \(M>0\)，根位于
\(0<\beta_c<\pi/2\)，所以端点附加质量降低临界转速。只将 \(m\) 换成
\(m+M\) 而保留原正切函数不正确，因为质量分布和端点边界不同。

在二次变分中，端点另贡献
\(-M\omega^2[\delta r(l_0)]^2\)，从同一混合边界得到相同临界根。图 3-2
的长度曲线采用
\(M=0,r_0=0\)；加入端点质量后需按新公式绘制，不能只把原曲线的横轴比例缩放。

\subsection{六、关键洞察}\label{ux516dux5173ux952eux6d1eux5bdf-2}

这里的 \(x\)
必须是自然状态材料坐标，不能直接拿转动后的空间坐标当质量坐标。这样
\(dm=(m/l_0)dx\) 始终成立，局部胡克定律也能正确写出。

连续介质题的三个层次是：材料坐标标记质量、形函数描述空间位置、截面内力连接相邻微元。把三个量混成同一个
\(x\)，几乎一定会写错线密度或伸长率。

\textbf{局部劲度为何与长度成反比。} 整根自然长度 \(l_0\) 的弹簧劲度为
\(k\)，截取自然长度 \(dx\) 的一小段，相同材料和截面积下其劲度为
\(kl_0/dx\)。因此微段伸长 \(dr-dx\) 时，张力为

\[
F=\frac{kl_0}{dx}(dr-dx)
=kl_0\left(\frac{dr}{dx}-1\right).
\]

若直接写
\(F=k(dr-dx)\)，等于把每一个微段都误当成与整根弹簧具有同一劲度，连续极限会得到趋于零的总伸长。

\textbf{内力梯度的方向。}
外侧微元所需离心力更大，内侧截面要拉住它外面的所有质量，所以张力由自由端向固定端增大。若材料坐标从内向外增加，则
\(dF/dx<0\)。在写微元方程以前先作这个定性判断，可以检查差分
\(F(x+dx)-F(x)\) 的符号。

\textbf{临界条件还包含张力非负。} 数学表达式在 \(\cos\beta\ne0\)
时形式上存在，但物理弹簧只能承受拉伸模型下的非负张力。基本支应满足
\(0\le\beta<\pi/2\)，此时形函数单调且
\(F(x)\ge0\)。越过首个奇点后出现的其他正切分支会伴随形函数振荡或局部压缩，已不属于题设的径向拉伸稳态。

\subsection{七、易错点}\label{ux4e03ux6613ux9519ux70b9-2}

\begin{itemize}
\tightlist
\item
  把整根弹簧张力误认为处处相等。
\item
  用变形后的线密度却没有同步修改微元长度。
\item
  忘记自由端张力为零对应 \(dr/dx=1\)，而不是 \(dr/dx=0\)。
\item
  只写出含正切的形式，却不讨论 \(\beta=\pi/2\) 的发散。
\item
  把整体劲度系数 \(k\) 直接当成任意微段的劲度系数。
\item
  忘记固定端距转轴为 \(r_0\)，把边界写成 \(r(0)=0\)。
\item
  只写 \(\cos\beta\ne0\)，却接受越过首个奇点后的非物理解分支。
\item
  没有先规定 \(x\) 从固定端指向自由端，导致 \(dF/dx\) 的符号反转。
\item
  把总长度误写成 \(r(1)\)，忘记减去固定端半径 \(r_0\)。
\end{itemize}

\subsection{八、课堂讲解建议}\label{ux516bux8bfeux5802ux8bb2ux89e3ux5efaux8bae-2}

先用一串相同小弹簧与小质点的离散模型建立直觉，再取连续极限。强调固定端给位置条件、自由端给力条件，两个边界条件性质不同。

解出结果后立即做 \(\omega\to0\)、\(r_0\to0\) 和 \(\beta\to\pi/2\)
三个极限检查，学生会看到公式中每一项分别承担什么物理作用。

\textbf{课堂推导顺序。}
先用三段离散弹簧说明内侧张力大于外侧张力，再引入自然坐标
\(x\)。在板书左侧固定写 \(dm=(m/l_0)dx\)，右侧写 \(dr\)
是变形后长度，提醒两者属于不同构形。完成 \(F=kl_0(dr/dx-1)\)
后，再对微元列力平衡，微分方程就会自然出现。

\textbf{边界条件的意义。} 固定端条件是位移型边界
\(r(0)=r_0\)；自由端条件是力型边界
\(F(l_0)=0\)。这与弦、杆、梁等连续体系的边值问题结构相同。学生如果把自由端写成
\(r(l_0)=l_0\)，说明仍把自然长度和空间位置混在一起。

\textbf{临界性的讲法。} \(\beta=\pi/2\)
不是普通共振频率，而是静态旋转平衡解失去有限性的点。可把
\(l/l_0=\tan\beta/\beta\)
曲线画出，让学生观察低速平缓、高速急剧上升的行为，并讨论真实材料为何会更早失效。

\textbf{离散到连续的课堂桥梁。} 先画三段弹簧，每段自然长度
\(l_0/3\)、质量 \(m/3\)、劲度
\(3k\)。从最外段开始列平衡，学生能直接看到内侧张力等于外侧所有微元离心力之和。再把段数改为
\(n\)，写出

\[
F_{j+1}-F_j=-\frac mn\omega^2r_j,
\]

最后才令
\(n\to\infty\)。这样微分方程是离散受力的自然极限，而不是突然出现的数学技巧。

\textbf{边界条件板书。}
固定端写``给位置''：\(r(0)=r_0\)；自由端写``不给位置、给力''：\(F(l_0)=0\)。随后由局部胡克定律把第二条翻译为
\(r'(l_0)=1\)。建议保留这一步翻译，不能直接报导数边界，否则学生不清楚数字
1 来自``零伸长率''而不是某个坐标归一化。

\textbf{四级核验。} 先验量纲：\(\beta\)
必须无量纲；再验低速：\(l\to l_0\)；再验无偏置：\(r_0\to0\)；最后验临界：\(\beta\to\pi/2^-\)
时长度与固定端张力同时发散。四级核验覆盖参数、边界与物理极限，适合训练连续介质解题习惯。

\subsection{九、变式训练}\label{ux4e5dux53d8ux5f0fux8badux7ec3-2}

\begin{enumerate}
\def\labelenumi{\arabic{enumi}.}
\tightlist
\item
  若外端再连一个质量 \(M\)，写出新的自由端边界条件。
\item
  对 \(\omega\sqrt{m/k}\ll1\) 作二阶展开，求长度的近似增量。
\item
  若弹簧线密度沿材料坐标线性变化，写出相应的变系数微分方程。
\end{enumerate}

\Needspace{3\baselineskip}
\textbf{变式解析要点}

\begin{enumerate}
\def\labelenumi{\arabic{enumi}.}
\tightlist
\item
  外端质量需要向心力，故自由端张力不再为零，而是
  \(F(l_0)=M\omega^2r(l_0)\)；将其代入局部胡克定律可得到新的混合边界条件。
\item
  使用 \(\tan\beta/\beta=1+\beta^2/3+O(\beta^4)\) 与
  \(\sec\beta-1=\beta^2/2+O(\beta^4)\)，可得
  \(\Delta l=(m\omega^2/k)(l_0/3+r_0/2)\)。
\item
  若 \(dm=\mu(x)dx\)，微元平衡变为
  \(dF/dx=-\mu(x)\omega^2r\)；局部本构若仍均匀，则最终得到变系数或非齐次边值问题，求解时仍须保留材料坐标。
\end{enumerate}

\newpage

\section{第 4 题}\label{ux7b2c-4-ux9898}

半径 \(r=1\,\mathrm m\) 的均匀电阻圆环固定在水平面内，单位长度电阻
\(\lambda=1\,\Omega/\mathrm m\)。一根单位长度电阻同为 \(\lambda\)
的均匀直棒以圆环直径为长度，其中心与环心重合，可绕过中心的竖直轴无摩擦转动；转动惯量为
\(mr^2\)，\(m=1\,\mathrm{kg}\)。圆环四等分点和转轴均接地。空间有竖直向下的匀强磁场
\(B=1\,\mathrm T\)。初始棒平分相邻接地点，角速度
\(\omega_0=10\,\mathrm{rad/s}\)。

\begin{enumerate}
\def\labelenumi{\arabic{enumi}.}
\tightlist
\item
  求初始流过直棒的电流 \(I_0\)。
\item
  足够长时间后，直棒停在标示的 \(A\) 区还是 \(B\) 区？
\end{enumerate}

\begin{figure}[H]
\centering
\includegraphics[width=0.58\linewidth,height=\textheight,keepaspectratio,alt={图 4-1　接地电阻圆环与可转动导电直棒。}]{../figures/question-04.pdf}
\caption*{图 4-1　接地电阻圆环与可转动导电直棒。}
\end{figure}

\subsection{一、题型定位}\label{ux4e00ux9898ux578bux5b9aux4f4d-3}

先把一侧半棒和两条接地圆弧画成等效电路，再求电流和力矩。用方位角积分阻尼，可确定多圈转动后的极限停角。

\subsection{二、知识点}\label{ux4e8cux77e5ux8bc6ux70b9-3}

\begin{itemize}
\tightlist
\item
  转动导体棒的运动电动势。
\item
  圆弧电阻的串并联等效。
\item
  安培力矩与楞次定律。
\item
  变量分离微分方程。
\item
  周期位置下的累计衰减。
\end{itemize}

\subsection{三、物理图像}\label{ux4e09ux7269ux7406ux56feux50cf-3}

棒的两半分别与相邻接地点构成闭合回路。棒转动切割磁感线产生电流，电流在磁场中受力，力矩恒阻碍转动。接地点把圆环分成四段，因此阻尼在每个象限重复。

\subsection{四、解题入口}\label{ux56dbux89e3ux9898ux5165ux53e3-3}

令棒相对某一象限平分线的角位移为 \(\theta\)，象限半角
\(\theta_0=\pi/4\)。先求一侧端点到地之间的等效电阻
\(R(\theta)\)，再求电流与力矩。

\subsection{五、解析}\label{ux4e94ux89e3ux6790-3}

\begin{figure}[H]
\centering
\includegraphics[width=0.578\linewidth,height=\textheight,keepaspectratio,alt={图 4-2　棒转过角度 \textbackslash theta 时的环路、电流与磁力矩。}]{../figures/solution-04.pdf}
\caption*{图 4-2　棒转过角度 \(\theta\) 时的环路、电流与磁力矩。}
\end{figure}

图 4-2
画出一侧半棒的等效电路：半棒电阻与两段接地圆弧的并联电阻串联；另一侧完全相同。下面先按这三条支路求一侧总电阻，再沿半棒从
\(0\) 到 \(r\) 积分动生电动势和阻力矩。

\Needspace{3\baselineskip}
\textbf{1. 一侧环路的等效电阻}

棒从中心到圆环的半根电阻为
\(\lambda r\)。棒端到相邻两个接地点的圆弧长度分别为
\(r(\theta_0-\theta)\)、\(r(\theta_0+\theta)\)，两段圆弧并联：

\[
R_{\rm arc}
=\lambda r\frac{(\theta_0-\theta)(\theta_0+\theta)}{2\theta_0}
=\lambda r\frac{\theta_0^2-\theta^2}{2\theta_0}.
\]

所以一侧总电阻

\[
\boxed{
R(\theta)=\lambda r
\frac{2\theta_0+\theta_0^2-\theta^2}{2\theta_0}
}. \tag{1}
\]

\Needspace{3\baselineskip}
\textbf{2. 转棒的运动电动势}

半根棒上距轴 \(\rho\) 处线元速度为 \(\omega\rho\)。运动电动势

\[
\mathcal E
=\int_0^rB\omega\rho,d\rho
=\frac12B\omega r^2.
\]

因而一侧电流

\[
I(\theta)=\frac{\mathcal E}{R(\theta)}
=\frac{Br\theta_0\omega}
{\lambda(2\theta_0+\theta_0^2-\theta^2)}. \tag{2}
\]

初始在象限平分线，\(\theta=0\)。代入
\(r=B=\lambda=1\)、\(\theta_0=\pi/4\)、\(\omega_0=10\)：

\[
\begin{aligned}
I_0
&=\frac{Br\theta_0\omega_0}{\lambda(2\theta_0+\theta_0^2)}\\
&=\frac{40}{8+\pi}\,\mathrm A\\
&\approx3.59\,\mathrm A.
\end{aligned}
\]

故

\[\boxed{I_0=40/(8+\pi)\,\mathrm A\approx3.59\,\mathrm A}.\]

\Needspace{3\baselineskip}
\textbf{3. 电磁阻尼方程}

半根棒的安培力矩为

\[
\tau_{1}=\int_0^rBI\rho,d\rho=\frac12BIr^2.
\]

两半棒力矩同向，总力矩大小

\[
\tau=BIr^2.
\]

由楞次定律，方向与 \(\omega\) 相反。转动方程

\[
mr^2\dot\omega=-BIr^2.
\]

代入 (2)，并用 \(\dot\theta=\omega\)：

\[
\dot\omega
=-\frac{rB^2\theta_0}{m\lambda}
\frac{\omega}{2\theta_0+\theta_0^2-\theta^2}.
\]

又因为

\[\dot\omega=\frac{d\omega}{d\theta}\dot\theta
=\omega\frac{d\omega}{d\theta},\]

约去非零的 \(\omega\)，得到

\[
d\omega=-\frac{rB^2\theta_0}{m\lambda}
\frac{d\theta}{2\theta_0+\theta_0^2-\theta^2}. \tag{3}
\]

\Needspace{3\baselineskip}
\textbf{4. 每跨过一个象限的角速度损失}

从一个平分线到下一个平分线，\(\theta\) 由 \(-\theta_0\) 变到
\(+\theta_0\)。积分 (3)：

\[
\Delta\omega
=-\frac{rB^2\theta_0}{m\lambda}
\int_{-\theta_0}^{\theta_0}
\frac{d\theta}{2\theta_0+\theta_0^2-\theta^2}.
\]

令 \(a^2=2\theta_0+\theta_0^2\)，使用

\[\int\frac{d\theta}{a^2-\theta^2}
=\frac{1}{2a}\ln\frac{a+\theta}{a-\theta},\]

可得

\[
\boxed{
\Delta\omega
=-\frac{rB^2\theta_0}{m\lambda a}
\ln\frac{a+\theta_0}{a-\theta_0}
}.
\]

代入数值：

\[\Delta\omega\approx-0.628\,\mathrm{rad/s}.\]

初始角速度可跨过的完整象限数约为

\[
N=\frac{\omega_0}{|\Delta\omega|}
\approx15.9.
\]

也就是说转棒越过 15 个象限后，在第 16 个象限内停止。按图中 \(A,B\)
区交替排列，第 16 个象限属于 \(B\) 区，因此

\[\boxed{\text{最终停在 }B\text{ 区}}.\]

\Needspace{3\baselineskip}
\textbf{能量方向检查}

电磁力矩的瞬时功率为

\[\tau\omega=-2I^2R(\theta)<0,\]

负号表示转动机械能转化为焦耳热。因此 \(\omega\)
必须单调减小，若积分结果给出正的
\(\Delta\omega\)，说明楞次定律方向写反。

\Needspace{3\baselineskip}
\textbf{初始电流的电路拆分}

当 \(\theta=0\) 时，两段接地圆弧长度都为 \(r\theta_0\)，并联电阻为

\[R_{\rm arc}(0)=\frac{\lambda r\theta_0}{2}.\]

串联半根棒电阻后

\[R(0)=\lambda r\left(1+\frac{\theta_0}{2}\right).\]

而半棒电动势是 \(Br^2\omega_0/2\)，所以

\[
I_0=\frac{Br^2\omega_0/2}
{\lambda r(1+\theta_0/2)},
\]

与正文的 \(Br\theta_0\omega_0/[\lambda(2\theta_0+\theta_0^2)]\)
完全等价。

\Needspace{3\baselineskip}
\textbf{阻尼积分的逐步化简}

令

\[a=\sqrt{\theta_0^2+2\theta_0}.\]

则

\[
\begin{aligned}
\int_{-\theta_0}^{\theta_0}\frac{d\theta}{a^2-\theta^2}
&=\left.\frac1{2a}\ln\frac{a+\theta}{a-\theta}
\right|_{-\theta_0}^{\theta_0}\\
&=\frac1a\ln\frac{a+\theta_0}{a-\theta_0}\\
&=\frac2a\ln\sqrt{\frac{a+\theta_0}{a-\theta_0}}.
\end{aligned}
\]

若在定义``半个象限''或``完整象限''时改变积分区间，前面的 2
会随之改变；因此区域计数必须与积分区间保持同一约定。

\Needspace{3\baselineskip}
\textbf{为什么每个象限损失相同角速度}

电阻函数只依赖
\(\theta^2\)，每经过一个接地点后电路拓扑重复，局部坐标重新从
\(-\theta_0\) 计到 \(+\theta_0\)。式 (3) 中 \(d\omega/d\theta\) 不显含
\(\omega\)，所以每个完整象限的 \(\Delta\omega\)
相同，而每个象限损失的转动能并不相同。

\Needspace{3\baselineskip}
\textbf{停止位置的定量判断}

跨过 15 个完整象限后剩余角速度

\[
\omega_{15}=\omega_0-15|\Delta\omega|
\approx10-15\times0.628
\approx0.58\,\mathrm{rad/s}.
\]

它小于跨过第 16 个完整象限所需的 \(0.628\,\mathrm{rad/s}\)，所以棒在第
16 象限内部停止，而不是恰停在接地点。

\Needspace{3\baselineskip}
\textbf{机械功率与焦耳热}

两个半棒回路的焦耳功率为 \(P_J=2I^2R(\theta)\)。磁力矩功率

\[P_{\rm mech}=-BIr^2\omega.\]

利用 \(I=(Br^2\omega/2)/R\)，立即得到

\[P_{\rm mech}=-2I^2R=-P_J.\]

转动机械能的减少与焦耳热完全相等。

\Needspace{3\baselineskip}
\textbf{全过程的能量积分}

若棒最终静止，且系统没有其他机械能输入，则从初态到终态积分上式：

\[
\int_0^{t_{\rm stop}}P_J\,dt
=-\int_0^{t_{\rm stop}}P_{\rm mech}\,dt
=\frac12J\omega_0^2.
\]

本题
\(J=mr^2=1\,\mathrm{kg,m^2}\)、\(\omega_0=10\,\mathrm{rad/s}\)，所以总焦耳热为

\[
Q_J=\frac12\times1\times10^2=50\,\mathrm J.
\]

这个结果与棒停在哪一个区域无关：区域只决定瞬时等效电阻和减速快慢，而总机械能损失由初、末状态决定。静止后
\(\omega=0\)，动生电动势、电流和磁阻力矩同时归零；停止角由此前累计的角速度损失确定，并不是某一区域存在额外的势能最低点。

\Needspace{3\baselineskip}
\textbf{极限停角与初始半个区间}

完整接地间隔损失
\(0.6282341553\,\mathrm{rad/s}\)。初始位于间隔中点，第一次到接地点只损失其一半；不能把第一段误当完整接地间隔。按``中点到中点''计，经过十五段后余量为
\(0.5764876705\)，还会过下一接地点，然后在 \(B\) 区趋于停止。

取初始棒轴为方位 \(\psi=0\)，逆时针为正，解最后一段积分可得

\[\psi_\infty=8\pi-0.1436800370\simeq24.98906119\,\mathrm{rad}.\]

因此最终棒轴接近水平，落在 \(B\) 区。图 4-1
把接地点画在四个对角方向，\(A\) 表示上下扇区、\(B\)
表示左右扇区；原稿把接地点直接标在 \(A,B\) 中心，破坏了初态几何。

在极限角附近 \(R(\theta)\)
有限且正，\(\dot\omega=-c_*\omega\)，所以速度随时间指数趋零，而不是在有限时刻突然变零。总焦耳热仍为
\(50\,\mathrm J\)，时间积分上限应理解为 \(\infty\)。

\Needspace{3\baselineskip}
\textbf{忽略棒电阻的变式存在不可积阻尼}

令 \(q=\pi/4\)，若半棒电阻取零，一侧只剩两圆弧并联，

\[R_{\rm arc}=\lambda r\frac{q^2-\theta^2}{2q}.\]

电流、力矩与角速度仍满足
\(I=B\omega r^2/(2R_{\rm arc})\)、\(mr^2\dot\omega=-BIr^2\)。于是

\[\frac{d\omega}{d\theta}
=-\frac{B^2rq}{m\lambda(q^2-\theta^2)}.\]

从某初角 \(\theta_i\) 积分到 \(\theta\)，角速度为

\[\omega(\theta)=\omega_i-\frac{B^2r}{2m\lambda}
\ln\frac{(q+\theta)(q-\theta_i)}{(q-\theta)(q+\theta_i)}.\]

在趋近接地点 \(\theta\to q^-\)
时，对数发散。因此有限初速度会在接地点之前出现 \(\omega=0\)
的极限角，不能继续假定以非零速度跨过该接地点，再把发散电流说成短暂经过。

对固定有限初速度，极限角严格小于
\(q\)，在那里圆弧电阻仍正，实际时间仍为指数趋停。若棒电阻或接触电阻为正，分母不再归零，每个接地间隔的角速度损失有限，恢复正文的多圈计数。这个极限说明``很小电阻''与``严格零电阻''对跨越接地点不是可随意互换的近似。

\subsection{六、关键洞察}\label{ux516dux5173ux952eux6d1eux5bdf-3}

棒端连接到两个相邻接地点的圆弧是并联，不是串联。阻尼力矩与角速度成正比，但比例系数随角位置周期改变；因此要按象限积分，而不能直接套指数衰减。

变量替换 \(\dot\omega=\omega,d\omega/d\theta\) 后，方程中的 \(\omega\)
恰好消去，这是``每跨过相同角区间损失相同角速度''的数学根源。

\textbf{电路拓扑必须先于动力学。}
一侧半棒从转轴到圆环接触点形成电源与棒电阻，接触点到相邻两个接地点的两段圆弧接在同一对节点之间，故为并联。另一侧半棒构成完全对称的第二支路。只有把节点画清，才能判断初始总电流、每侧电流以及磁力矩中的因子
2；如果先套动生电动势公式，往往会在电阻串并联上失分。

\textbf{阻尼为何不是时间指数。} 在某一固定角度，确有
\(\tau=-C(\theta)\omega\)；但棒转动时 \(C(\theta)\) 随位置变化，所以
\(J\dot\omega=-C(\theta)\omega\) 不是常系数方程。改用
\(\dot\omega=(d\omega/d\theta)\dot\theta\) 后，\(\dot\theta=\omega\)
与右侧的 \(\omega\)
消去，得到每个相同角区间有相同的角速度损失。若研究的是随时间变化，\(\omega(t)\)
仍不是简单指数。

\textbf{有限停止的含义。} 因为每跨过一个完整象限损失固定的
\(|\Delta\omega|\)，棒在有限个象限后无法越过下一个接地点区间，并在区间内部渐近静止，有限时间内角速度仍为正。停止后动生电动势为零，电流和磁力矩也为零，因此最终角度不是静电或磁势能极小点，只是耗散历程决定的位置。

\subsection{七、易错点}\label{ux4e03ux6613ux9519ux70b9-3}

\begin{itemize}
\tightlist
\item
  忘记半根棒自身的电阻 \(\lambda r\)。
\item
  只算一半棒的磁力矩。
\item
  把 \(d\omega/dt\) 直接换成 \(d\omega/d\theta\)，漏掉因子
  \(\dot\theta=\omega\)。
\item
  用 \(15.9\) 直接四舍五入决定区域，应该看已经完整跨过多少象限。
\item
  忘记四个接地点使电路拓扑每 \(\pi/2\) 重复一次。
\item
  把半棒电动势 \(Br^2\omega/2\) 写成 \(Br^2\omega\)。
\item
  把圆弧的弧长 \(r\theta\) 误写成角度 \(\theta\)，造成电阻量纲错误。
\item
  把两侧半棒的电流先相加后再计算每侧力矩，重复计数。
\item
  由 \(N=15.9\) 断言``转过 15.9
  个象限''，忽略最后不足一个象限的积分应从边界重新计算。
\end{itemize}

\subsection{八、课堂讲解建议}\label{ux516bux8bfeux5802ux8bb2ux89e3ux5efaux8bae-3}

先在图上用两种颜色标出一侧棒电阻与两条并联圆弧，再写电路。第二问用``每一格损失相同的角速度，而非相同能量''帮助学生判断跨区次数。

可让学生先用数值积分画出 \(\omega\)
随方位角的折线，观察每个象限下降量重复，再回到解析积分解释原因。

\textbf{课堂推导顺序。}
先只画棒的一半：棒电阻与两段圆弧并联后的电阻串联。求出一侧电流后，再借对称性复制到另一半。磁力矩也先积分半棒，再乘
2。这样每一个``2''的来源都能在图上找到，不会混在最终公式中。

\textbf{位置变量的约定。} 每进入一个新象限，都把局部坐标重新定义为
\(-\theta_0\le\theta\le\theta_0\)。必须在图上标出局部零点是象限平分线，而非某个固定接地点。否则学生会误以为
\(R(\theta)\) 在多圈运动中不断变化而不重复。

\textbf{区域判定。} 得到 \(N=15.9\) 后，先写``完整通过 15 个象限，在第
16 个内停止''，再结合初始方向沿图逐格编号。不要直接把 15.9 四舍五入成
16，因为停止在第 16 区内部与通过第 16 区是不同事件。

\textbf{半棒电动势的图上推导。} 在离轴距离 \(x\) 处取微元，速度
\(v_x=\omega x\)，动生电动势

\[
d\mathcal E=B\omega x\,dx.
\]

从轴到端积分得 \(\mathcal E_{1/2}=B\omega r^2/2\)。建议让学生在图 4-2
上同时标出 \(x\)、\(d\bm l\)、\(\bm v\) 和 \(\bm B\)，再用
\(\bm v\times\bm B\)
判断哪一端电势高。这样电动势大小和电流方向来自同一幅图，而不是两套互不相干的记忆。

\textbf{磁力矩的逐元计算。} 同一微元受力 \(dF=IB\,dx\)，对轴力矩
\(d\tau=xIB\,dx\)，故一侧半棒

\[
\tau_{1/2}=\int_0^r xIB\,dx=\frac12IBr^2.
\]

两侧相加得
\(\tau=IBr^2\)，方向由楞次定律判为阻碍转动。课堂上把电动势积分与力矩积分上下对齐，学生会看到两个
\(1/2\) 分别来自同一个 \(\int_0^r x\,dx\)，但总力矩还要乘两侧。

\textbf{能量核验。} 由 \(I=\mathcal E_{1/2}/R\) 和两支路焦耳热
\(2I^2R\)，可验证
\(-\tau\omega=2I^2R\)。这个恒等式同时检查电流、力矩和因子
2；若两边不相等，说明电路或力矩计数必有错误。

\subsection{九、变式训练}\label{ux4e5dux53d8ux5f0fux8badux7ec3-3}

\begin{enumerate}
\def\labelenumi{\arabic{enumi}.}
\tightlist
\item
  若圆环只有上下两个点接地，重新写 \(R(\theta)\)。
\item
  若忽略棒电阻，讨论接近接地点时电流是否发散。
\item
  若磁场只覆盖圆环的上半部分，判断阻尼周期如何改变。
\end{enumerate}

\Needspace{3\baselineskip}
\textbf{变式解析要点}

\begin{enumerate}
\def\labelenumi{\arabic{enumi}.}
\tightlist
\item
  只剩两个接地点后，棒端到接地点的两条弧仍并联，但拓扑周期由 \(\pi/2\)
  改为 \(\pi\)；应重新定义局部角坐标而不是直接沿用原积分上下限。
\item
  忽略棒电阻时，接触点靠近接地点会使一条圆弧电阻趋于零。理想模型给出很大电流，但同时作用角区间趋窄；需检查积分是否可积，并说明真实接触电阻会消除发散。
\item
  磁场只覆盖上半圆后，只有位于磁区内的棒段产生电动势和受磁力，左右或上下支路不再完全对称。必须按棒与磁区的交长度分段积分，阻尼周期也随覆盖区域改变。
\end{enumerate}

\newpage

\section{第 5 题}\label{ux7b2c-5-ux9898}

一线圈密绕成细圆环，圆环半径 \(a=0.1\,\mathrm m\)，线圈截面积
\(S=2\,\mathrm{cm^2}\)，匝数
\(N=200\)，内部为空气。电流在一小段时间内以恒定速率

\[\frac{dI}{dt}=10\,\mathrm{A/s}\]

变化。同时在细圆环的几何中心放置一个初始静止的质子。忽略量子效应与相对论效应，求此时质子的瞬时加速度。

\begin{figure}[H]
\centering
\includegraphics[width=0.84\linewidth,height=\textheight,keepaspectratio,alt={图 5-1　细圆环线圈的空间结构与位于几何中心的质子。}]{../figures/question-05.pdf}
\caption*{图 5-1　细圆环线圈的空间结构与位于几何中心的质子。}
\end{figure}

\subsection{一、题型定位}\label{ux4e00ux9898ux578bux5b9aux4f4d-4}

逐匝计算中心矢势时，各匝磁矩沿大圆环切向，指向中心的位矢沿径向；两者叉乘同向叠加，时变部分给轴向感生电场。

\subsection{二、知识点}\label{ux4e8cux77e5ux8bc6ux70b9-4}

\begin{itemize}
\tightlist
\item
  单匝线圈的磁偶极矩与磁矢势。
\item
  环形分布矢量的方向判断与叠加。
\item
  \(\bm E=-\partial\bm A/\partial t\)。
\item
  轴对称条件下中心点感生电场的方向。
\item
  电磁类比与互易性。
\item
  带电粒子瞬时加速度。
\end{itemize}

\subsection{三、物理图像}\label{ux4e09ux7269ux7406ux56feux50cf-4}

大圆环上的每一匝小线圈都包围截面积 \(S\)。第 \(i\) 匝的磁偶极矩
\(\bm m_i\) 沿大圆环的切向，而从该匝中心指向几何中心的位矢 \(\bm R_i\)
沿径向。由于 \(\bm m_i\perp\bm R_i\)，且绕环一周时
\(\bm m_i\times\bm R_i\)
始终沿同一条对称轴，各匝在中心产生的磁矢势同向相加。电流变化后，轴向磁矢势随时间变化，于是几何中心出现轴向感生电场。质子初始静止，瞬时磁力为零，只有电场力。

\subsection{四、解题入口}\label{ux56dbux89e3ux9898ux5165ux53e3-4}

不要把整个细圆环当成一个磁偶极矩为 \(NIS\) 的平面线圈。正确入口是任选第
\(i\) 匝，写出它在几何中心的磁矢势 \(\bm A_i\)；再用圆环对称性证明所有
\(\bm A_i\) 同向，因而可以直接乘以匝数 \(N\)。

\subsection{五、解析}\label{ux4e94ux89e3ux6790-4}

\begin{figure}[H]
\centering
\includegraphics[width=0.67\linewidth,height=\textheight,keepaspectratio,alt={图 5-2　单匝磁偶极矩、中心位矢与轴向感生电场的几何关系。}]{../figures/solution-05.pdf}
\caption*{图 5-2　单匝磁偶极矩、中心位矢与轴向感生电场的几何关系。}
\end{figure}

图 5-2 中每一小匝的磁矩 \(\bm m_i\) 沿大圆环切向，而指向几何中心的位矢
\(\bm R_i\) 沿径向，因此 \(\bm m_i\times\bm R_i\)
对所有匝都沿同一轴向。下面先按这组矢量方向叠加磁矢势，再对时间求导。

\Needspace{3\baselineskip}
\textbf{1. 建立单匝线圈的矢量关系}

记大圆环半径为 \(a\)。在第 \(i\) 匝小线圈处建立柱坐标单位矢量
\(\hat{\bm r}_i\)、\(\hat{\bm\varphi}_i\)、\(\hat{\bm z}\)。选定电流绕向后，第
\(i\) 匝的磁偶极矩可写成

\[
\bm m_i=IS\hat{\bm\varphi}_i. \tag{1}
\]

从该匝中心 \(C_i\) 指向几何中心 \(P\) 的位矢为

\[
\bm R_i=-a\hat{\bm r}_i,
\qquad |\bm R_i|=a. \tag{2}
\]

一个磁偶极在场点产生的磁矢势是

\[
\bm A_i(P)
=\frac{\mu_0}{4\pi}
\frac{\bm m_i\times\bm R_i}{|\bm R_i|^3}. \tag{3}
\]

把式 (1)、(2) 代入，并使用柱坐标叉乘关系

\[
\hat{\bm\varphi}_i\times\hat{\bm r}_i=-\hat{\bm z},
\]

得到

\[
\begin{aligned}
\bm m_i\times\bm R_i
&=IS\hat{\bm\varphi}_i\times
(-a\hat{\bm r}_i)\\
&=ISa\hat{\bm z},
\end{aligned}
\]

从而

\[
\bm A_i(P)=\frac{\mu_0IS}{4\pi a^2}\hat{\bm z}. \tag{4}
\]

这里最关键的是：虽然 \(\hat{\bm r}_i\) 与 \(\hat{\bm\varphi}_i\)
随匝的位置改变，但式 (4) 的方向对每一匝都相同。于是 \(N\)
匝的磁矢势可直接相加：

\[
\bm A(P)
=\sum_{i=1}^{N}\bm A_i(P)
=\frac{\mu_0NIS}{4\pi a^2}\hat{\bm z}. \tag{5}
\]

\Needspace{3\baselineskip}
\textbf{2. 由时变矢势求感生电场}

中心点处的感生电场为

\[
\bm E(P)=-\frac{\partial\bm A(P)}{\partial t}.
\]

对式 (5) 求时间偏导：

\[
\bm E(P)
=-\frac{\mu_0NS}{4\pi a^2}
\frac{dI}{dt}\hat{\bm z}. \tag{6}
\]

因此电场大小为

\[
E
=\frac{\mu_0NS}{4\pi a^2}
\left|\frac{dI}{dt}\right|.
\]

代入

\[
\mu_0=4\pi\times10^{-7}\,\mathrm{H/m},
\]

\[N=200,\qquad S=2\times10^{-4}\,\mathrm{m^2},\]

\[a=0.1\,\mathrm m,\qquad \frac{dI}{dt}=10\,\mathrm{A/s},\]

得到

\[
\begin{aligned}
E
&=\frac{4\pi\times10^{-7}\times200\times2\times10^{-4}\times10}
{4\pi\times(0.1)^2}\\
&=4.0\times10^{-6}\,\mathrm{V/m}.
\end{aligned}
\]

\Needspace{3\baselineskip}
\textbf{3. 质子加速度}

质子初始速度为零，故

\[\bm F=q_p\bm E+q_p\bm v\times\bm B=q_p\bm E.\]

取

\[q_p=e=1.602\times10^{-19}\,\mathrm C,\]

\[m_p=1.673\times10^{-27}\,\mathrm{kg},\]

则

\[
\begin{aligned}
a_p
&=\frac{eE}{m_p}\\
&=\frac{1.602\times10^{-19}\times4.0\times10^{-6}}
{1.673\times10^{-27}}\\
&\approx3.83\times10^2\,\mathrm{m/s^2}.
\end{aligned}
\]

所以

\[\boxed{a_p\approx3.83\times10^2\,\mathrm{m/s^2}},\]

方向沿细圆环的对称轴，具体指向由线圈绕向和 \(dI/dt\) 的正负共同决定。

\Needspace{3\baselineskip}
\textbf{量纲检查}

式 (6) 中 \(\mu_0NS(dI/dt)/a^2\) 的量纲为

\[
\frac{\mathrm{T\,m}}{\mathrm A}
\frac{\mathrm{m^2\,A/s}}{\mathrm{m^2}}
=\mathrm{T\,m/s}=\mathrm{V/m},
\]

与电场量纲一致。

\Needspace{3\baselineskip}
\textbf{方向的确定}

先由单匝电流方向用右手定则确定 \(\bm m_i\)，再由
\(\bm m_i\times\bm R_i\) 确定 \(\bm A_i\) 的轴向。因为
\(\bm E=-\partial\bm A/\partial t\)，当 \(dI/dt>0\) 时，\(\bm E\) 与
\(\bm A\) 反向；当 \(dI/dt\)
变号时，电场和质子加速度同时反向，但大小不变。

\Needspace{3\baselineskip}
\textbf{从磁偶极矢势到法拉第定律}

第 \(i\) 匝的磁偶极矢势满足

\[
\bm A(\bm r)=\frac{\mu_0}{4\pi}
\frac{\bm\mu\times\bm r}{r^3}.
\]

在中心点，\(\bm m_i\perp\bm R_i\)，所以
\(A_i=\mu_0IS/(4\pi a^2)\)。对任意闭合曲线，

\[
\oint\bm E\cdot d\bm l
=-\frac{d}{dt}\oint\bm A\cdot d\bm l
=-\frac{d}{dt}\int\bm B\cdot d\bm S,
\]

因此矢势方法与法拉第定律完全等价，只是更适合本题的空间几何。

\Needspace{3\baselineskip}
\textbf{近似适用范围}

磁偶极近似要求每一匝小线圈的线度远小于大圆环半径
\(a\)；准静态近似还要求电流变化的时间尺度远大于
\(a/c\)，从而可以忽略传播延迟与辐射场。题目所说``细圆环''和``一小段时间内变化率保持不变''分别对应这两个近似。

\Needspace{3\baselineskip}
\textbf{数值代入的逐级检查}

先算所有小线圈面积与电流变化率的乘积：

\[
NS\frac{dI}{dt}
=200\times2\times10^{-4}\times10
=0.4\,\mathrm{A\,m^2/s}.
\]

再算电场：

\[
E=\frac{10^{-7}}{(0.1)^2}\times0.4
=4.0\times10^{-6}\,\mathrm{V/m}.
\]

最后乘质子的荷质比 \(e/m_p\approx9.58\times10^7\,\mathrm{C/kg}\)，得到
\(a_p\approx383\,\mathrm{m/s^2}\)。分三步代入能减少指数错误。

\Needspace{3\baselineskip}
\textbf{有限变化时间内的速度增量}

若电流变化率在 \(0\le t\le\Delta t\)
内保持不变，且质子位移小到可以继续使用中心点电场，则加速度也是常量：

\[
a_p=\frac{e\mu_0NS}{4\pi m_pa^2}
\left|\frac{dI}{dt}\right|.
\]

因此质子在这段时间内获得的速度增量为

\[
\begin{aligned}
\Delta v
&=a_p\Delta t\\
&=\frac{e\mu_0NS}{4\pi m_pa^2}
\left|\frac{dI}{dt}\right|\Delta t\\
&=\frac{e\mu_0NS}{4\pi m_pa^2}|\Delta I|.
\end{aligned}
\]

这个形式说明：在准静态且位移可忽略的条件下，最终速度增量只由总电流改变量
\(\Delta I\)
决定，与电流变化得快或慢无关；变化越快，瞬时加速度越大，但作用时间相应越短。质子增加的动能来自维持线圈电流变化的外部电源所做的功，并非由静磁场直接做功。

\Needspace{3\baselineskip}
\textbf{中心点近似的自洽条件}

取质子初始位置为 \(z=0\)。在 \(0\le t\le\Delta t\)
内把中心电场视为常量时，轴向位移为

\[
|\Delta z|=\frac12a_p(\Delta t)^2.
\]

为了仍可使用中心点场值，需要 \(|\Delta z|\ll a\)，即

\[
\Delta t\ll\sqrt{\frac{2a}{a_p}}.
\]

代入 \(a=0.1\,\mathrm m\)、\(a_p=3.83\times10^2\,\mathrm{m/s^2}\)：

\[
\sqrt{\frac{2a}{a_p}}
=\sqrt{\frac{0.2}{3.83\times10^2}}
\approx2.3\times10^{-2}\,\mathrm s.
\]

所以当电流按给定变化率维持的时间远小于约 \(2.3\times10^{-2}\,\mathrm s\)
时，既可忽略质子离开中心后场的空间变化，也能保持非相对论近似。若作用时间更长，就必须把
\(E(z)\) 的位置依赖写入运动方程，而不能继续使用匀加速公式。

从做功关系看，在这段短时间内

\[
\Delta K=\int q_p\bm E\cdot d\bm r,
\]

能量由驱动线圈电流变化的外部电源经感生电场传递给质子。磁场力始终与速度垂直，本身不改变质子的动能。

\Needspace{3\baselineskip}
\textbf{中心场方向与所采用的电场模型}

若图 5-2 取小匝磁矩为 \(+\hat\varphi\)，中心矢势沿
\(+\hat z\)；\(dI/dt>0\) 时感生电场必须沿 \(-\hat z\)。原图向上的 \(E\)
与文字叉乘相反，现分开画 \(A\) 向上、\(E\) 向下。

完整电场为
\(\bm E=-\partial_t\bm A-\nabla\phi\)。本题计算理想均匀绕组的感生场，并忽略供电导线、绕组表面电荷和外加静电场；仅给
\(dI/dt\) 不能指定这些额外静电贡献。所谓中心加速度
\(383\,\mathrm{m/s^2}\)
应按这个理想感生场模型理解，而不能说任何实际供电结构都严格相同。

匝截面为 \(2\,\mathrm{cm^2}\)，等效小圆半径
\(\sqrt{S/\pi}\simeq8.0\,\mathrm{mm}\)，与 \(a=0.1\,\mathrm m\) 的比约
\(0.08\)，支持小匝偶极近似。反接半数匝时中心感生场相消，这是逐匝方向计算的直接检验。

\subsection{六、关键洞察}\label{ux516dux5173ux952eux6d1eux5bdf-4}

本题不是把整个装置压缩成一个总磁偶极，而是叠加 \(N\)
个方向随位置改变的小磁偶极的矢势。磁偶极矩本身绕环一周后矢量和为零，但中心点的
\(\bm A_i\propto\bm m_i\times\bm R_i\)
却全部同向；``磁矩和为零''并不意味着``中心矢势为零''。

本题给的是电流变化率而不是电流本身，所以质子加速度由 \(d\bm A/dt\)
决定，与该时刻 \(I\)
的绝对值无关。恒定电流只产生静态场，对初始静止质子没有瞬时电场力。

\textbf{为什么总磁矩为零仍有中心矢势。} 环上各匝的 \(\bm m_i\)
沿切向分布，绕一周的矢量和确实为零；但场点矢势不是只看
\(\sum\bm m_i\)，而是看

\[
\sum_i\frac{\bm m_i\times\bm R_i}{R_i^3}.
\]

在中心点所有 \(R_i=a\)，并且 \(\bm m_i\times\bm R_i\)
同向。因此不能先把磁矩相加再把整个装置当成一个点偶极。这个例子非常适合说明``多极矩的总和''与``有限距离处的场''不是同一层次的对象。

\textbf{矢势不是可随意省略的中间量。}
本题中心附近的磁场结构并不容易用一个简单高斯面或安培环路得到，而
\(\bm A_i\) 的方向叠加却非常直接。选取对称规范后，中心轴向的
\(\bm E=-\partial\bm A/\partial t\)
具有明确物理意义；最终可由法拉第定律闭合积分核验，因此并非依赖规范的可观测量。

\textbf{瞬时加速度与后续轨迹分开。} 题目只问初始瞬时加速度，此时
\(\bm v=0\)，磁力严格为零。质子开始运动后，若空间同时存在磁场，则
\(q\bm v\times\bm B\)
会逐渐出现，电场也可能随位置改变。不能把初始匀加速结果无限延长到任意时间。

\subsection{七、易错点}\label{ux4e03ux6613ux9519ux70b9-4}

\begin{itemize}
\tightlist
\item
  把整个环形线圈误当成所有磁偶极矩方向相同的平面线圈。
\item
  直接相加 \(\bm m_i\)，没有相加真正决定中心场的
  \(\bm m_i\times\bm R_i\)。
\item
  把磁偶极矢势的 \(1/a^2\) 写成磁场的 \(1/a^3\)。
\item
  质子初始静止仍写磁力 \(qvB\)。
\item
  只给加速度大小，不说明轴向方向随线圈绕向或 \(dI/dt\) 符号反转。
\item
  使用厘米平方时忘记 \(1\,\mathrm{cm^2}=10^{-4}\,\mathrm{m^2}\)。
\item
  看到``圆环中心''便误套普通圆电流线圈中心磁场 \(B=\mu_0NI/(2a)\)。
\item
  把磁矢势方向画成圆环平面内切向，与三维结构不符。
\item
  用 \(E\propto1/a^2\) 的中心结果讨论远离中心的任意位置。
\end{itemize}

\subsection{八、课堂讲解建议}\label{ux516bux8bfeux5802ux8bb2ux89e3ux5efaux8bae-4}

先拿出一匝，分别画出 \(\bm m_i\)、\(\bm R_i\) 与
\(\bm m_i\times\bm R_i\)；然后把这组三矢量沿大圆环移动一周，让学生观察叉乘结果始终沿同一轴向。之后再从
\(\bm E=-\partial\bm A/\partial t\)
说明电流变化率为何直接决定轴向电场。强调质子初始静止，所以无需计算磁力。

可将结果与重力加速度比较：\(383\,\mathrm{m/s^2}\approx39g\)。数值很大但作用时间很短，有助于学生区分``瞬时加速度大''和``最终速度一定大''。

\textbf{课堂推导顺序。} 先写单匝磁矩
\(\bm m_i=IS\hat{\bm\varphi}_i\)，再写中心位矢
\(\bm R_i=-a\hat{\bm r}_i\)，完成一次叉乘以后才乘
\(N\)，最后对时间求导。每一步只引入一个新物理量。数值代入时分成
\(NS(dI/dt)\)、\(E\)、\(a_p\) 三行，避免 \(N,S,a\)
的指数全部挤在一个长式中。

\textbf{法拉第定律的联系。} 说明 \(-\partial\bm A/\partial t\) 与
\(-d\Phi/dt\)
并不是两套定律。此处若直接寻找穿过某个简单平面的总磁通，几何关系很复杂；矢势方法允许先逐匝计算中心点的贡献，再利用对称性相加，更适合本题。

\textbf{几何方向核对。} 在大圆环相对的两侧，\(\bm m_i\) 与 \(\bm R_i\)
都同时反向，所以叉乘 \(\bm m_i\times\bm R_i\)
不变号。这是判断各匝贡献相加而非抵消的最快方法。

\textbf{三维图的讲法。}
先让学生辨认大圆环的平面和每一匝小线圈的平面：小线圈的法向沿大圆环切线，因此
\(\bm m_i\) 并不沿总轴线。随后在图 5-2 只挑一匝画 \(\bm m_i\)
与指向中心的 \(\bm R_i\)，用右手定则得到轴向
\(\bm A_i\)。再移动到对径点，两矢量同时反向而叉乘不变。两点核对后即可用连续对称性推广到整环。

\textbf{数值计算的三道闸门。} 第一闸门是面积换算
\(2\,\mathrm{cm^2}=2\times10^{-4}\,\mathrm{m^2}\)；第二闸门是
\(\mu_0/(4\pi)=10^{-7}\)；第三闸门是质子荷质比
\(e/m_p\approx9.58\times10^7\,\mathrm{C/kg}\)。分别计算

\[
NS\frac{dI}{dt}=0.4\,\mathrm{A\,m^2/s},
\qquad E=4.0\times10^{-6}\,\mathrm{V/m},
\]

再得到 \(a_p\approx383\,\mathrm{m/s^2}\)。分级代入能快速定位数量级错误。

\textbf{可观测方向。} 若把所有绕组方向反接，\(\bm m_i\) 与 \(\bm A\)
同时反向；若保持绕向而令 \(dI/dt\) 变号，\(\bm E\)
反向。题目没有给绕线手性时，只能确定加速度大小和``沿对称轴''，不应擅自指定向上或向下。

\Needspace{22\baselineskip}

\subsection{九、变式训练}\label{ux4e5dux53d8ux5f0fux8badux7ec3-4}

\begin{enumerate}
\def\labelenumi{\arabic{enumi}.}
\tightlist
\item
  将线圈匝数加倍、每匝截面积减半，判断中心质子加速度是否改变。
\item
  把一半匝数的绕向反接，求中心点的瞬时电场。
\item
  若换成电子，分别比较加速度大小与方向。
\end{enumerate}

\Needspace{3\baselineskip}
\textbf{变式解析要点}

\begin{enumerate}
\def\labelenumi{\arabic{enumi}.}
\tightlist
\item
  中心电场只依赖乘积 \(NS\)，所以 \(N\to2N\)、\(S\to S/2\)
  后不变；但若改变截面尺寸破坏``小线圈线度远小于
  \(a\)''的条件，还要重新评估偶极近似。
\item
  一半绕组反接后，两组 \(\bm A_i\)
  等大反向，理想对称情形中心矢势及其时间导数均为零，质子初始加速度为零。
\item
  电子电荷为负，受力方向与质子相反；加速度大小之比为
  \(m_p/m_e\approx1836\)。应强调这是相同电场下的瞬时比较。
\end{enumerate}

\newpage

\section{第 6 题}\label{ux7b2c-6-ux9898}

激光陀螺利用沿闭合光路传播的两束反向激光测量相对于惯性空间的转动。先考虑中心到边的距离为
\(r_0\)、周长为 \(L\)、面积为 \(A=r_0L/2\)
的正方形环路，环路绕垂直于平面的轴以角速度 \(\omega\) 逆时针转动。

\begin{enumerate}
\def\labelenumi{\arabic{enumi}.}
\tightlist
\item
  求逆时针光束绕一圈的光程 \(L_+\)。
\item
  求顺、逆两束光的光程差，并在 \(r_0\omega\ll c\) 下保留一阶项。
\item
  对
  \(\omega=15^\circ/\mathrm h\)、\(\lambda=0.70\,\mu\mathrm m\)、\(A=600\times300\,\mathrm{m^2}\)，说明光程差约为
  \(\lambda/4\)。
\item
  改成闭合环形共振腔后，求 \(L_+-L_-\)。
\item
  用 \(A,\lambda,L,\omega\) 表示两束光的拍频 \(f_--f_+\)。
\end{enumerate}

\begin{figure}[H]
\centering
\includegraphics[width=0.523\linewidth,height=\textheight,keepaspectratio,alt={图 6-1　随平台转动的正方形闭合光路。}]{../figures/question-06.pdf}
\caption*{图 6-1　随平台转动的正方形闭合光路。}
\end{figure}

\subsection{一、题型定位}\label{ux4e00ux9898ux578bux5b9aux4f4d-5}

光传播时镜面也在转动，一阶时间修正沿闭合路径累积成面积项。先求传播差，再把它换成环形腔同一纵模的频率差。

\subsection{二、知识点}\label{ux4e8cux77e5ux8bc6ux70b9-5}

\begin{itemize}
\tightlist
\item
  运动边界追及关系。
\item
  小参数一阶近似。
\item
  萨格纳光程差。
\item
  环形腔共振条件。
\item
  双频拍频。
\end{itemize}

\subsection{三、物理图像}\label{ux4e09ux7269ux7406ux56feux50cf-5}

逆时针光束与环路同向，必须多走一点才能回到移动的分束点；顺时针光束与分束点迎面相遇，路程略短。这个差异在闭合光路四边累加，只与光路包围的有向面积有关。

\subsection{四、解题入口}\label{ux56dbux89e3ux9898ux5165ux53e3-5}

先用一条边说明追赶关系，再推广到四边。因为题目只需一阶精度，可把边界切向速度在各边的投影保留到
\(\omega r_0\) 阶。

\subsection{五、解析}\label{ux4e94ux89e3ux6790-5}

\begin{figure}[H]
\centering
\includegraphics[width=0.66\linewidth,height=\textheight,keepaspectratio,alt={图 6-2　光束追赶运动边界所产生的传播时间差。}]{../figures/solution-06.pdf}
\caption*{图 6-2　光束追赶运动边界所产生的传播时间差。}
\end{figure}

图 6-2
分别画出逆时针光束追赶分束点和顺时针光束迎向分束点的传播关系。后文的
\(L_+\)、\(L_-\) 将按图中边界速度取相反符号，二者相减得到萨格纳光程差。

\Needspace{3\baselineskip}
\textbf{1. 逆时针光束的光程}

设逆时针光束绕一圈传播时间为 \(t_+\)，光实际传播距离为

\[L_+=ct_+.\]

在这段时间内，分束点沿环路方向前进。对正方形逐边投影并求和，边界沿闭合光路方向的等效速度为
\(2A\omega/L\)，所以追赶关系为

\[
L_+=L+\frac{2A\omega}{L}\frac{L_+}{c}.
\]

整理为

\[
L_+\left(1-\frac{2A\omega}{Lc}\right)=L,
\]

所以

\[
\boxed{
L_+=\frac{L}{1-2A\omega/(Lc)}
}. \tag{1}
\]

对于正方形，\(L=8r_0\)、\(A=4r_0^2=r_0L/2\)，故

\[
L_+=\frac{L}{1-r_0\omega/c}.
\]

在 \(r_0\omega/c\ll1\) 时，利用

\[\frac1{1-x}=1+x+O(x^2),\]

得到

\[
L_+\approx L+\frac{2A\omega}{c}. \tag{2}
\]

\Needspace{3\baselineskip}
\textbf{2. 顺时针光束与光程差}

顺时针光束迎向分束点，符号反向：

\[
L_-=\frac{L}{1+2A\omega/(Lc)}
\approx L-\frac{2A\omega}{c}. \tag{3}
\]

因此一阶光程差

\[
\begin{aligned}
\Delta L
&=L_+-L_-\\
&\approx\left(L+\frac{2A\omega}{c}\right)
-\left(L-\frac{2A\omega}{c}\right)\\
&=\frac{4A\omega}{c}.
\end{aligned}
\]

故

\[\boxed{\Delta L=4A\omega/c}.\]

\Needspace{3\baselineskip}
\textbf{3. 地球自转数值估算}

先把角速度换成国际单位：

\[
\omega=15^\circ/\mathrm h
=\frac{15\pi/180}{3600}\,\mathrm{rad/s}
=7.272\times10^{-5}\,\mathrm{rad/s}.
\]

面积

\[A=600\times300=1.8\times10^5\,\mathrm{m^2}.\]

于是

\[
\begin{aligned}
\Delta L
&=\frac{4\times1.8\times10^5\times7.272\times10^{-5}}
{2.998\times10^8}\\
&\approx1.75\times10^{-7}\,\mathrm m.
\end{aligned}
\]

而

\[\frac\lambda4=\frac{0.70\times10^{-6}}4
=1.75\times10^{-7}\,\mathrm m.\]

故

\[\boxed{\Delta L\approx\lambda/4}.\]

\Needspace{3\baselineskip}
\textbf{4. 环形共振腔}

闭合共振腔的几何追赶关系不变，所以在上述等效追赶模型中，光程写为

\[
L_\pm=\frac{L}{1\mp2A\omega/(Lc)}.
\]

在等效常投影速度模型内，两式代数相减：

\[
\begin{aligned}
L_+-L_-
&=L\left[
\frac1{1-x}-\frac1{1+x}
\right]\\
&=L\frac{2x}{1-x^2},
\end{aligned}
\]

其中

\[x=\frac{2A\omega}{Lc}.\]

所以

\[
\boxed{
L_+-L_-=
\frac{4A\omega/c}{1-[2A\omega/(Lc)]^2}
\approx\frac{4A\omega}{c}
}.
\]

\Needspace{3\baselineskip}
\textbf{5. 共振频率与拍频}

共振条件

\[L_+=m\lambda_+,\qquad L_-=m\lambda_-.\]

频率

\[f_+=\frac{c}{\lambda_+}=\frac{mc}{L_+},\qquad
f_-=\frac{mc}{L_-}.\]

故

\[
f_--f_+
=mc\left(\frac1{L_-}-\frac1{L_+}\right).
\]

由 \(m\approx L/\lambda\)，且

\[
\frac1{L_-}-\frac1{L_+}
=\frac{4A\omega}{L^2c},
\]

得到

\[
\boxed{
f_--f_+=\frac{4A\omega}{\lambda L}
}.
\]

\Needspace{3\baselineskip}
\textbf{给定数值的拍频核验}

当 \(L=0.40\,\mathrm m\)、正方形面积
\(A=L^2/16=0.010\,\mathrm{m^2}\)、\(\lambda=6328\,\text{\AA}=6.328\times10^{-7}\,\mathrm m\)
时，

\[
\Delta f=\frac{4A\omega}{\lambda L}
\approx11.5\,\mathrm{Hz}.
\]

若实际有效面积或入射纬度投影不同，应把 \(A\)
替换成有效矢量面积；这正是装置标定时必须说明的几何因素。

\Needspace{3\baselineskip}
\textbf{一阶近似的误差数量级}

小参数是

\[\epsilon=\frac{2A\omega}{Lc}.\]

地球自转实验中 \(\epsilon\) 通常远小于 \(10^{-10}\)。略去项从
\(\epsilon^2\) 开始，因此用

\[L_\pm\approx L(1\pm\epsilon)\]

计算光程差远比装置自身的机械、温度误差精确。物理上的难点不在近似，而在保持腔长与有效面积稳定。

\Needspace{3\baselineskip}
\textbf{逐边推导面积因子}

对第 \(i\) 条直边，端点位置为
\(\bm r_i,\bm r_{i+1}\)。转动速度在边方向上的平均投影给出一阶传播时间修正

\[
\delta t_i=\frac{\bm\omega\cdot
(\bm r_i\times\bm r_{i+1})}{c^2}.
\]

沿闭合多边形求和：

\[
\sum_i\bm r_i\times\bm r_{i+1}=2\bm A.
\]

顺、逆两束符号相反，因此总时间差

\[\Delta t=\frac{4\bm\omega\cdot\bm A}{c^2},\]

乘以 \(c\) 即得
\(\Delta L=4\bm\omega\cdot\bm A/c\)。这也证明一阶结果与光路的具体形状无关。

\Needspace{3\baselineskip}
\textbf{从光程差到相位差}

普通干涉仪的相位差

\[\Delta\phi=\frac{2\pi}{\lambda}\Delta L
=\frac{8\pi A\omega}{\lambda c}.\]

当 \(\Delta L=\lambda/4\)
时，\(\Delta\phi=\pi/2\)，正好对应四分之一个条纹周期。

\Needspace{3\baselineskip}
\textbf{等效追赶模型的展开}

令 \(x=2A\omega/(Lc)\)，则

\[
L_+-L_-=\frac{L}{1-x}-\frac{L}{1+x}
=\frac{2Lx}{1-x^2}.
\]

因为

\[\frac1{1-x^2}=1+x^2+O(x^4),\]

所以

\[
\Delta L=\frac{4A\omega}{c}
\left[1+\left(\frac{2A\omega}{Lc}\right)^2+\cdots\right].
\]

常投影速度模型的展开下一项为三阶；实际运动镜面的高阶光路仍须另算。一阶面积公式要求整个光程内边界位移足够小。

\Needspace{3\baselineskip}
\textbf{拍频的方向信息}

逆时针转动时 \(L_+>L_-\)，同一纵模序号下 \(\lambda_+>\lambda_-\)，因此
\(f_+<f_-\)。反转平台转向后，拍频绝对值不变、符号反转，故激光陀螺既能测转速也能判定转向。

\Needspace{3\baselineskip}
\textbf{正方形追赶式只保证一阶精度}

把各边速度投影统一代成常量
\(2A\omega/L\)，可得到正文的分式；但运动镜面在传播期间还会改变方向，因此它不是旋转正方形光路的严格高阶解。正文分母
\(1-\epsilon^2\)
只是这一等效追赶模型的代数结果，不能用它声称真实装置三阶修正已被精确求出。

可靠的一阶结论是

\[L_\pm=L\pm\frac{2\bm A\cdot\bm\omega}{c}+O(L\epsilon^2),\quad
\Delta L=\frac{4\bm A\cdot\bm\omega}{c}\quad\text{（一阶）}.\]

逆时针平台在顶边的速度指向左，原解图右指的速度画成了相反转向。题图也将两束光改画成完整的相反环路，并用颜色配
\(L_+,L_-\)，不能把两束都沿同一转向画在不同半边。

\subsection{六、关键洞察}\label{ux516dux5173ux952eux6d1eux5bdf-5}

普通干涉仪测的是一次传播的微小相位差；环形激光器把同一几何效应转化成两组腔模的稳定频差，因而可持续积累和高精度测量。

萨格纳效应的一阶结果只依赖有向面积而不依赖光路形状。真正参与测量的是
\(\bm A\cdot\bm\omega\)，因此装置姿态也决定灵敏度。

\textbf{光速不变与追赶关系并不冲突。} 在惯性系中光始终以 \(c\)
传播，镜面和分束点在光传播期间发生位移，所以闭合条件是``光走过的路程等于初始几何周长再加或减边界位移''。式中的
\(c\mp\omega r_0\)
是解追赶时间得到的有效分母，不是说光相对局部镜面的基本速度变成了
\(c\mp v\)。

\textbf{一阶面积定律的来源。} 对第 \(i\)
条边，端点在传播时间内的位移投影产生修正
\(\delta L_i\propto\bm\omega\cdot(\bm r_i\times\bm r_{i+1})/c\)。沿闭合多边形求和后

\[
\frac12\sum_i\bm r_i\times\bm r_{i+1}=\bm A,
\]

故顺逆光程差只保留
\(4\bm A\cdot\bm\omega/c\)。边数、边长分配和具体形状都被闭合求和吸收到有向面积中。

\textbf{普通干涉与环形腔的读出量不同。} 普通装置把 \(\Delta L\)
转为相位差
\(2\pi\Delta L/\lambda\)，观察条纹；环形腔要求两个方向各自满足纵模条件，转动导致共振频率分裂，最后读出拍频。物理源头相同，但测量量从``单次相位''变成``持续频差''，灵敏度与稳定性也随之改变。

\subsection{七、易错点}\label{ux4e03ux6613ux9519ux70b9-5}

\begin{itemize}
\tightlist
\item
  把 \(c\) 理解为相对移动镜面的速度并做伽利略相加。
\item
  光程差漏掉顺、逆两方向的双倍差，写成 \(2A\omega/c\)。
\item
  角速度从度每小时换成弧度每秒时漏除 \(3600\)。
\item
  拍频结果中漏掉周长 \(L\)。
\item
  把 \(15^\circ/\mathrm h\) 直接当成 \(15\,\mathrm{rad/s}\)。
\item
  忘记共振腔中两个方向的纵模序号在近似处理中取相同整数 \(m\)。
\item
  把面积 \(A\) 与边长或周长混为同一量纲，导致 \(4A\omega/c\) 量纲错误。
\item
  在计算拍频时把自由空间波长与腔内光程条件重复乘折射率。
\item
  忽略有向面积，认为转轴反向后光程差大小和符号都不变。
\end{itemize}

\subsection{八、课堂讲解建议}\label{ux516bux8bfeux5802ux8bb2ux89e3ux5efaux8bae-5}

先用``跑道上的追赶者''建立单边直觉，再用面积公式概括四边。数值题应完整写单位转换，让学生看到巨大面积为何仍只产生亚微米光程差。

讲到环形激光器时把``光程差---波长差---频率差''画成一条因果链，强调后者可以用电子计数精确测量，因而比直接观察微小条纹移动更实用。

\textbf{课堂推导顺序。}
先只求传播时间差，不急于谈干涉。用运动分束点的追赶关系得到 \(L_\pm\)
后，再依次乘 \(2\pi/\lambda\) 得相位差、套 \(L_\pm=m\lambda_\pm\)
得波长差、最后用 \(f=c/\lambda\) 得拍频。三个层次应分开板书。

\textbf{面积公式的重点。} 对正方形可以逐边追赶，但最终应指出
\(\sum\bm r_i\times\bm r_{i+1}=2\bm A\)。这一步把特例提升为一般闭合光路，并说明为什么弯曲或改变边数不会改变一阶结果，只要有向面积不变。

\textbf{实验量级。} 地球自转角速度只有
\(7.27\times10^{-5}\,\mathrm{rad/s}\)，普通干涉仪需要巨大面积才能得到可见条纹；环形激光器测频差，允许用较小面积获得可计数信号。应让学生区分``基本物理效应大小''和``测量方案灵敏度''。

\textbf{近似展开逐行写。} 令

\[
\varepsilon=\frac{2A\omega}{Lc}\ll1,
\]

则 \(L_+=L/(1-\varepsilon)\)、\(L_-=L/(1+\varepsilon)\)。课堂上明确使用

\[
\frac1{1\mp\varepsilon}=1\pm\varepsilon+O(\varepsilon^2),
\]

再相减得到
\(\Delta L=2L\varepsilon=4A\omega/c\)。这样学生能看到忽略的是二阶项，也能理解为何两个方向的一阶修正相加而不是抵消。

\textbf{单位换算板书。} 地球角速度应先写

\[
15^\circ/\mathrm h
=15\frac{\pi}{180}\frac1{3600}\,\mathrm{rad/s},
\]

面积写成 \(1.8\times10^5\,\mathrm{m^2}\)，波长写成
\(0.70\times10^{-6}\,\mathrm m\)。每个量单独换算后再代入，能避免把度、小时和微米同时塞进一行造成数量级错误。

\textbf{拍频的符号。} 若规定 \(\omega>0\)
为逆时针，哪个方向频率更高取决于 \(L_\pm\) 的定义。实际测量常关注
\(|f_--f_+|\)；理论推导则应保留符号，用它判断转动方向。课堂上应区分``频差大小''与``有向频差''。

\textbf{从光程到条纹移动的完整链条。} 两束光重新叠加时的相位差为

\[
\Delta\phi=\frac{2\pi}{\lambda}\Delta L
=\frac{8\pi A\omega}{\lambda c}.
\]

对应的条纹移动数为

\[
\Delta N=\frac{\Delta\phi}{2\pi}
=\frac{4A\omega}{\lambda c}.
\]

所以题目给出的 \(\Delta L\approx\lambda/4\) 等价于条纹移动约 \(1/4\)
条。这里的``光程差''和``条纹数''相差一个波长因子，不能把 \(\lambda/4\)
直接说成四分之一弧度。

\textbf{共振腔频差逐步推导。} 定义小量

\[
\varepsilon=\frac{2A\omega}{Lc},
\]

并取同一纵模序号 \(m\)。两方向的有效光程与共振波长满足

\[
L_+=\frac{L}{1-\varepsilon}=m\lambda_+,
\qquad
L_-=\frac{L}{1+\varepsilon}=m\lambda_-.
\]

因此

\[
\begin{aligned}
f_+&=\frac{c}{\lambda_+}
=\frac{mc}{L_+}
=\frac{mc}{L}(1-\varepsilon),\\
f_-&=\frac{c}{\lambda_-}
=\frac{mc}{L_-}
=\frac{mc}{L}(1+\varepsilon).
\end{aligned}
\]

相减得

\[
\begin{aligned}
f_--f_+
&=\frac{2mc}{L}\varepsilon\\
&=\frac{2c}{\lambda}\frac{2A\omega}{Lc}\\
&=\boxed{\frac{4A\omega}{\lambda L}}.
\end{aligned}
\]

第二行使用静止时 \(m=L/\lambda\)。这个逐步推导能明确周长 \(L\)
为什么出现在拍频分母中，也说明环形腔测得的是同一几何光程差造成的模频率分裂。

\textbf{适用近似的层次。} 基本的一阶萨格纳式要求平台边缘速度远小于
\(c\)，并忽略光在介质中的色散、镜面有限尺寸和结构弹性变形。若环路中有静止介质或介质随平台共同转动，还需明确采用相位速度还是群速度及拖曳效应。本题为空气或真空近似，使用
\(c\) 即可，不应把工程修正提前混入基础模型。

\subsection{九、变式训练}\label{ux4e5dux53d8ux5f0fux8badux7ec3-5}

\begin{enumerate}
\def\labelenumi{\arabic{enumi}.}
\tightlist
\item
  对任意平面闭合光路证明一阶结果仍为 \(4\bm A\cdot\bm\omega/c\)。
\item
  若转轴与光路法线夹角为 \(\alpha\)，写出拍频。
\item
  两个面积相同但周长不同的共振腔，比较其拍频灵敏度。
\end{enumerate}

\Needspace{3\baselineskip}
\textbf{变式解析要点}

\begin{enumerate}
\def\labelenumi{\arabic{enumi}.}
\tightlist
\item
  将光路离散为多边形，逐边写运动端点的传播时间修正，再用鞋带公式或叉乘和识别有向面积；曲线可由多边形极限得到。
\item
  只取角速度沿面积法向的分量，故
  \(\Delta f=4A\omega\cos\alpha/(\lambda L)\)；当光路平面与转轴平行时一阶信号为零。
\item
  在 \(A\)
  相同、波长相同时，\(\Delta f\propto A/L\)，周长较短者频差更大。但工程上还需兼顾镜面损耗、模宽和可实现几何，不能只用单一比例判断最终精度。
\end{enumerate}

\newpage

\section{第 7 题}\label{ux7b2c-7-ux9898}

把室内空调视为工作在室外温度 \(T_H\) 与室内温度 \(T_L\)
之间的理想可逆制冷机。每循环从室内吸热 \(Q_L\)，向室外放热
\(Q_H\)，电泵做功 \(W\)。空调输入功率为 \(P\)。

\begin{enumerate}
\def\labelenumi{\arabic{enumi}.}
\tightlist
\item
  用 \(T_H,T_L,P\) 表示最大室内热量移除速率。
\item
  若室外以速率 \(\dot Q_{\rm in}=k(T_H-T_L)\)
  向室内传热，求稳态最低室温。
\item
  \(k=173\,\mathrm{W/K}\)、室外 \(40^\circ\mathrm C\)、室内目标
  \(25^\circ\mathrm C\)，求所需最小功率。
\end{enumerate}

\begin{figure}[H]
\centering
\includegraphics[width=0.76\linewidth,height=\textheight,keepaspectratio,alt={图 7-1　空调从室内吸热并向室外放热的能量流。}]{../figures/question-07.pdf}
\caption*{图 7-1　空调从室内吸热并向室外放热的能量流。}
\end{figure}

\subsection{一、题型定位}\label{ux4e00ux9898ux578bux5b9aux4f4d-6}

制冷量和输入电功率通过卡诺COP联系，再与墙体漏热比较。稳态最低室温是两者相等的温度，所有温度比都使用开尔文。

\subsection{二、知识点}\label{ux4e8cux77e5ux8bc6ux70b9-6}

\begin{itemize}
\tightlist
\item
  可逆循环总熵变为零。
\item
  制冷机性能系数。
\item
  功率形式的能量守恒。
\item
  稳态热平衡。
\item
  温标必须使用开尔文。
\end{itemize}

\subsection{三、物理图像}\label{ux4e09ux7269ux7406ux56feux50cf-6}

温差越大，墙体向室内漏热越快；同时理想制冷机的性能系数越低。最低温度出现在两者刚好平衡处，再低一点就会出现漏热大于制冷量，无法维持。

\subsection{四、解题入口}\label{ux56dbux89e3ux9898ux5165ux53e3-6}

对一个可逆循环写高、低温热库的总熵变，再结合
\(W=Q_H-Q_L\)。将每循环量除以循环时间，即得到功率关系。

\subsection{五、解析}\label{ux4e94ux89e3ux6790-6}

\begin{figure}[H]
\centering
\includegraphics[width=0.818\linewidth,height=\textheight,keepaspectratio,alt={图 7-2　可逆制冷机的热流与功率关系。}]{../figures/solution-07.pdf}
\caption*{图 7-2　可逆制冷机的热流与功率关系。}
\end{figure}

图 7-2 规定低温端热流率 \(\dot Q_L\) 进入制冷机、输入功率 \(P=\dot W\)
加入循环、高温端热流率 \(\dot Q_H\) 排向室外。下面先对单个循环使用
\(Q_L\)、\(W\)、\(Q_H\)，再除以循环时间转化为图中的功率量，两个层次的符号一一对应。

\Needspace{3\baselineskip}
\textbf{1. 最大热量移除速率}

可逆循环中工作物质回到初态，熵变为零；两个热库的总熵变也为零：

\[
\frac{Q_H}{T_H}-\frac{Q_L}{T_L}=0.
\]

因此

\[
\frac{Q_H}{T_H}=\frac{Q_L}{T_L},
\]

即

\[Q_H=\frac{T_H}{T_L}Q_L.\]

一个循环的输入功

\[
W=Q_H-Q_L
=\left(\frac{T_H}{T_L}-1\right)Q_L
=\frac{T_H-T_L}{T_L}Q_L.
\]

所以制冷性能系数

\[
\mathrm{COP}=\frac{Q_L}{W}
=\frac{T_L}{T_H-T_L}.
\]

把能量除以时间，输入功率 \(P=\dot W\)，得到最大制冷功率

\[
\boxed{
\dot Q_{L,\max}
=\frac{T_L}{T_H-T_L}P
}.
\]

\Needspace{3\baselineskip}
\textbf{2. 最低稳态室温}

令

\[\Delta T=T_H-T_L>0.\]

墙体漏热速率

\[\dot Q_{\rm in}=k\Delta T.\]

在最低可维持温度的临界稳态，漏热与最大制冷量相等：

\[
k\Delta T
=\frac{T_L}{\Delta T}P
=\frac{T_H-\Delta T}{\Delta T}P.
\]

两边乘 \(\Delta T\)：

\[k\Delta T^2=P(T_H-\Delta T).\]

移到同一边：

\[
\Delta T^2+\frac Pk\Delta T-\frac PkT_H=0.
\]

取正根：

\[
\Delta T
=\frac{-P/k+\sqrt{(P/k)^2+4(P/k)T_H}}{2}.
\]

整理为

\[
\boxed{
\Delta T=\frac{P}{2k}
\left(\sqrt{1+\frac{4kT_H}{P}}-1\right)
}.
\]

所以最低室温

\[
\boxed{
T_{L,\min}=T_H-\frac{P}{2k}
\left(\sqrt{1+\frac{4kT_H}{P}}-1\right)
}.
\]

\Needspace{3\baselineskip}
\textbf{3. 数值计算}

室外温度

\[T_H=40+273=313\,\mathrm K,\]

目标室温

\[T_L=25+273=298\,\mathrm K,\]

温差

\[\Delta T=15\,\mathrm K.\]

由稳态方程

\[k\Delta T=\frac{T_L}{\Delta T}P,\]

解出

\[
P=k\frac{\Delta T^2}{T_L}.
\]

代入：

\[
\begin{aligned}
P
&=173\times\frac{15^2}{298}\,\mathrm W\\
&\approx130.6\,\mathrm W.
\end{aligned}
\]

故

\[\boxed{P_{\min}\approx131\,\mathrm W}.\]

\Needspace{3\baselineskip}
\textbf{极限检查}

当 \(P\to0\) 时，\(\Delta T\to0\)，室内不能低于室外；当 \(P\to\infty\)
时，在理想模型中可维持更大温差。真实空调还受换热器、压缩机与结霜等限制，实际功率必高于这里的可逆下限。

\Needspace{3\baselineskip}
\textbf{数值题的能流核验}

目标温差 \(15\,\mathrm K\) 时，墙体漏热为

\[\dot Q_{\rm in}=k\Delta T=173\times15=2595\,\mathrm W.\]

卡诺制冷系数

\[\mathrm{COP}=\frac{298}{15}\approx19.87.\]

所以最小输入功率

\[P=\frac{2595}{19.87}\approx130.6\,\mathrm W,\]

与二次方程反解完全一致。

\Needspace{3\baselineskip}
\textbf{最低温度根的选择}

二次方程另一个根为

\[
\Delta T_-=\frac{-P/k-\sqrt{(P/k)^2+4(P/k)T_H}}{2}<0.
\]

它意味着 \(T_L>T_H\)，与制冷情景的定义矛盾，必须舍去。保留的正根还应满足
\(0<\Delta T<T_H\)，这样绝对温度 \(T_L=T_H-\Delta T\) 才为正。

\Needspace{3\baselineskip}
\textbf{非可逆空调的修正}

若实际性能系数只有卡诺值的 \(\eta\) 倍，\(0<\eta<1\)，则

\[
\dot Q_L=\eta\frac{T_L}{T_H-T_L}P.
\]

稳态方程变为

\[k\Delta T^2=\eta P(T_H-\Delta T),\]

所需电功率是理想下限的 \(1/\eta\) 倍。

\Needspace{3\baselineskip}
\textbf{稳态曲线的唯一交点}

以 \(\Delta T\) 为横轴，漏热功率 \(k\Delta T\) 单调上升；最大制冷功率

\[P\frac{T_H-\Delta T}{\Delta T}\]

单调下降。两曲线在 \(0<\Delta T<T_H\)
内只有一个交点，因此最低稳态温度唯一。交点两侧的净热流方向都把系统推回交点，所以该稳态稳定。

\textbf{完整能量---熵联立。} 任意时间 \(dt\) 内有
\(dQ_H=dQ_L+Pdt\)，可逆熵关系为 \(dQ_H/T_H=dQ_L/T_L\)。消去 \(dQ_H\)
后得到

\[
\left(\frac{T_H}{T_L}-1\right)
\frac{dQ_L}{dt}=P,
\]

这就是正文第一问，且清楚地区分了热量和功。

\Needspace{3\baselineskip}
\textbf{动态稳定性与更准确的温标核算}

取 \(F(T_L)=k(T_H-T_L)-PT_L/(T_H-T_L)\)，房间热容为 \(C>0\) 时
\(C\dot T_L=F\)。导数

\[F'(T_L)=-k-\frac{PT_H}{(T_H-T_L)^2}<0,\]

所以稳态小温度扰动会衰减；若给定 \(C\)，线性恢复时间为
\(C/[k+PT_H/(T_H-T_L)^2]\)。不需假设热量能被``转成冷量''，三股热流的能量账目已足够。

使用 \(273.15\)
换算，\(T_H=313.15\,\mathrm K\)、\(T_L=298.15\,\mathrm K\)，最小功率
\(173\times225/298.15\simeq130.555\,\mathrm W\)，仍约
\(131\,\mathrm W\)。图 7-2
的虚线箭头表示室外向室内的漏热，稳态时其大小等于制冷量；把热漏关系代入可逆制冷功率，得到正文中室温的二次方程。

\subsection{六、关键洞察}\label{ux516dux5173ux952eux6d1eux5bdf-6}

题目给的是空调输入电功率，不是制冷功率。可逆极限下，空调从室内移走的热功率可以大于电功率，其比值就是制冷性能系数，并不违反能量守恒，因为同时向室外放出
\(\dot Q_H=\dot Q_L+P\)。

``最低温度''并不是空调瞬间能达到的温度，而是稳态方程的边界：该温度处最大制冷量恰好抵消漏热，再低就无法维持。

\textbf{三组方程分属三个层次。} \(\dot Q_H=\dot Q_L+P\)
是能量守恒；\(\dot Q_H/T_H=\dot Q_L/T_L\)
是可逆循环的熵平衡；\(\dot Q_L=k(T_H-T_L)\)
是房间稳态热平衡。三式不能混成一个``效率公式''直接背诵。先由前两式得到设备能力，再与第三式比较，模型层次才清楚。

\textbf{制冷系数为何可以大于 1。}
空调输入的功并不是被``转化成等量冷量''，而是驱动工质把室内热量搬到室外。因此

\[
\mathrm{COP}=\frac{\dot Q_L}{P}
=\frac{T_L}{T_H-T_L}
\]

可以远大于
1。真正受第二定律限制的是它不能超过卡诺值，而不是不能超过百分之百。

\textbf{稳态根的物理筛选。} 令 \(\Delta T=T_H-T_L>0\) 后，平衡式为
\(k\Delta T^2+P\Delta T-PT_H=0\)。代数上有两个根，但负根对应
\(\Delta T<0\)，即室内比室外更热，与``制冷最低温度''的分支不符。选择根必须结合定义域，而不是看到正负号就机械取加号。

\textbf{无量纲化显示控制参数。} 定义

\[
y=\frac{\Delta T}{T_H},
\qquad
p=\frac{P}{kT_H},
\]

则稳态方程化为

\[
y^2+py-p=0,
\]

物理解为

\[
y=\frac{\sqrt{p^2+4p}-p}{2}. \tag{A}
\]

所有设备功率、房间漏热与室外温度的影响都浓缩在无量纲比值 \(p\) 中。\(p\)
小表示设备能力相对漏热弱，\(p\)
大表示设备能力强。无量纲化不仅简化代数，也能直接比较不同房间和不同空调。

\textbf{弱功率与强功率极限。} 当 \(p\ll1\) 时，式 (A) 展开为

\[
y=\sqrt p-\frac p2+O(p^{3/2}),
\]

即

\[
\Delta T\approx\sqrt{\frac{PT_H}{k}}.
\]

最低温差随功率只按平方根增长，说明增加少量功率的收益会逐渐下降。当
\(p\gg1\) 时，

\[
y=1-\frac1p+O(p^{-2}),
\]

于是

\[
T_{L,\min}=T_H(1-y)\approx\frac{kT_H^2}{P}.
\]

即使功率很大，最低温度也只按 \(1/P\)
接近绝对零度，任何有限功率都不能达到 \(T_L=0\)。

\textbf{稳定性判据。} 若室温略高于平衡值，卡诺制冷量 \(PT_L/(T_H-T_L)\)
增大，而漏热 \(k(T_H-T_L)\)
减小，净制冷使温度下降；若室温略低，制冷能力减弱而漏热增强，温度回升。因此物理解对应稳定平衡。把代数根与微小扰动后的热流方向结合，能更准确地理解``可维持的最低温度''。

\subsection{七、易错点}\label{ux4e03ux6613ux9519ux70b9-6}

\begin{itemize}
\tightlist
\item
  把热机效率 \(1-T_L/T_H\) 直接套给制冷机。
\item
  温度用摄氏度代入 \(T_L/(T_H-T_L)\)。
\item
  第二问只写漏热 \(k\Delta T\)，不与制冷能力联立。
\item
  二次方程取到负温差根。
\item
  把墙体热导系数 \(k\) 与玻尔兹曼常量或热容混淆。
\item
  数值题用 \(T_L=25\) 代入 COP，而没有换成 \(298\,\mathrm K\)。
\item
  把 \(173\,\mathrm{W/^{\circ}C}\)
  当作房间热容，错误写成升降温速率方程。
\item
  把``最低温度''理解为有限时间内必然达到，而没有说明它是理想稳态极限。
\item
  计算 \(Q_H\) 时忘记加输入功率，写成 \(Q_H=Q_L\)。
\end{itemize}

\subsection{八、课堂讲解建议}\label{ux516bux8bfeux5802ux8bb2ux89e3ux5efaux8bae-6}

让学生先画三条能量流箭头并写
\(\dot Q_H=\dot Q_L+P\)，再谈熵。最后比较小温差和大温差时的
COP，解释为什么降低最后几度最困难。

数值题先算漏热 \(2595\,\mathrm W\)，再算 COP 约
19.9，最后相除得到输入功率。这样比直接套合并公式更能展示物理量之间的关系。

\textbf{课堂推导顺序。} 首先画出三股能量流并写
\(Q_H=Q_L+W\)；其次只写两个热库的熵变，得到
\(Q_H/T_H=Q_L/T_L\)；第三步消去 \(Q_H\) 得
COP。到第二问再引入墙体漏热，避免能量守恒、熵关系和房间热平衡三组方程同时出现。

\textbf{温标检查。} 可故意让学生比较用 \(25/15\) 与 \(298/15\) 得到的
COP，前者约 1.67、后者约
19.9。温差可以用摄氏度差，但绝对温度比必须用开尔文，这个数值对比能强化记忆。

\textbf{理想值的含义。} \(131\,\mathrm W\)
是可逆下限，不是常见空调的真实功率。应解释``最大可能效率''授权使用卡诺极限；若题目给实际性能系数，则必须以实际
COP 代替，不能继续把卡诺值当设备参数。

\textbf{能流图与公式一一对应。} 在图 7-2 中，左侧冷库流出
\(\dot Q_L\)，中间装置输入 \(P\)，右侧热库接收
\(\dot Q_H\)。要求学生沿箭头写

\[
\dot Q_H=\dot Q_L+P,
\]

再把两个热库的熵变分别写在箭头下方。图上每一支流都在方程中出现，能避免把
\(P\) 和 \(\dot Q_L\) 混为一谈。

\textbf{极限检查。} 当 \(T_L\to T_H\) 时，卡诺 COP
发散，表示搬运极小温差所需理想功趋小；当 \(T_L\to0\) 时，COP
趋零，接近绝对零度需要无限代价。对最低温度公式，\(P\to0\) 时
\(T_{L,\min}\to T_H\)，而 \(P\to\infty\)
时温差虽增加但不能由有限功率达到绝对零度。这些极限可用来检查二次方程根是否取对。

\textbf{数值题的解释顺序。} 先算房间在 \(15\,\mathrm K\) 温差下的漏热

\[
\dot Q_L=173\times15=2595\,\mathrm W,
\]

再算理想 COP 为 \(298/15\approx19.87\)，最后得到
\(P\approx130.6\,\mathrm W\)。最后补一句：实际空调存在不可逆性，所需电功率必大于此值。这样答案既有计算，也有现实含义。

\textbf{学生卡点诊断。} 若结果约为
\(1557\,\mathrm W\)，通常是误用了热机效率；若约为
\(2595\,\mathrm W\)，通常把制冷量直接当输入功率；若约为
\(10^2\,\mathrm W\)
但温标写成摄氏度，则只是偶然得到相近数量级，仍须纠正公式含义。

\textbf{课堂中的反事实检查。} 可问：若墙体完全绝热，即
\(k=0\)，稳态最低温度方程会怎样？理想模型给出只要持续工作便可不断降温，但越接近绝对零度
COP 越小、降温越慢。若输入功率为零，房间最终只能回到
\(T_H\)。这两个边界情形能帮助学生区分``设备循环极限''和``房间漏热平衡''两种限制。

\textbf{从稳态扩展到动态。} 若房间等效热容为 \(C\)，动态方程可写为

\[
C\frac{dT_L}{dt}
=k(T_H-T_L)-P\frac{T_L}{T_H-T_L}.
\]

稳态解就是右侧为零的温度。题目没有给
\(C\)，因此只能求稳态而不能求降温所需时间。此处适合提醒学生：出现
\(\mathrm{W/K}\) 的 \(k\) 是热导，不是热容；只有给出单位为
\(\mathrm{J/K}\) 的 \(C\) 才能讨论时间演化。

\textbf{备课中的现实修正。} 真实空调的 COP
低于卡诺值，并随室内外温差、换热器温度和压缩机工况变化。题目明确``达到最大可能效率''，才允许使用可逆卡诺模型。课堂最后说明这一点，能防止学生把
\(131\,\mathrm W\)
误认为实际设备选型结果，同时理解竞赛题是在求热力学下限。

\subsection{九、变式训练}\label{ux4e5dux53d8ux5f0fux8badux7ec3-6}

\begin{enumerate}
\def\labelenumi{\arabic{enumi}.}
\tightlist
\item
  若实际 COP 是卡诺值的 \(\eta\) 倍，改写最低温度方程。
\item
  若漏热规律为 \(k(T_H-T_L)+q_0\)，求稳态条件。
\item
  给定实际输入功率与实测 COP，反求房间等效热导系数。
\end{enumerate}

\Needspace{3\baselineskip}
\textbf{变式解析要点}

\begin{enumerate}
\def\labelenumi{\arabic{enumi}.}
\tightlist
\item
  写 \(\dot Q_L=\eta P T_L/(T_H-T_L)\)，再与 \(k(T_H-T_L)\)
  联立；\(\eta<1\) 会使可维持的最低温度升高。
\item
  稳态式改为
  \(PT_L/(T_H-T_L)=k(T_H-T_L)+q_0\)，乘去分母后仍是关于温差的二次方程，但常数项改变。
\item
  实测时 \(\dot Q_L=P\,\mathrm{COP}_{\rm real}\)，故
  \(k=P\,\mathrm{COP}_{\rm real}/(T_H-T_L)\)。若房间还有人员、电器等内热源，应先从总制冷量中扣除相应热功率。
\end{enumerate}

\newpage

\section{第 8 题}\label{ux7b2c-8-ux9898}

在惯性系 \(S\) 中，一长直圆筒半径为 \(R\)，固有表面电荷面密度为
\(\lambda_1/(2\pi R)\)（按扫描参考解口径；若题给的是 \(S\)
系观测密度，见解析中的另一口径），圆筒以速度 \(v\) 沿 \(+z\)
方向运动；中轴处有线电荷密度为 \(\lambda_2\) 的长直导线，在 \(S\)
中静止。点 \(P\) 距中轴 \(d\)。惯性系 \(S'\) 相对 \(S\) 以速度 \(u\) 沿
\(+z\) 方向运动。

\begin{enumerate}
\def\labelenumi{\arabic{enumi}.}
\tightlist
\item
  在 \(S\) 中分别求 \(d>R\) 与 \(d<R\) 时 \(P\) 点的电场和磁场。
\item
  求 \(S'\) 中圆筒、直导线的速度与线电荷密度。
\item
  当 \(d>R\) 时，求使 \(S'\) 中只有电场而无磁场的 \(u\) 及存在条件。
\item
  求使 \(S'\) 中只有磁场而无电场的 \(u\) 及存在条件。
\item
  是否存在某个 \(S'\) 使电场、磁场同时为零？
\end{enumerate}

\begin{figure}[H]
\centering
\includegraphics[width=0.701\linewidth,height=\textheight,keepaspectratio,alt={图 8-1　运动带电圆筒、静止轴线电荷与共轴运动参考系。}]{../figures/question-08.pdf}
\caption*{图 8-1　运动带电圆筒、静止轴线电荷与共轴运动参考系。}
\end{figure}

\subsection{一、题型定位}\label{ux4e00ux9898ux578bux5b9aux4f4d-7}

先明确线密度在哪个参考系给出，再把密度与电流作为一对量变换。筒内外分别计算，纯电或纯磁参考系由总线流的类型决定。

\subsection{二、知识点}\label{ux4e8cux77e5ux8bc6ux70b9-7}

\begin{itemize}
\tightlist
\item
  高斯定律求柱对称电场。
\item
  安培环路定律求轴向电流磁场。
\item
  相对论速度合成。
\item
  长度收缩与线电荷密度变换。
\item
  电磁场不变量。
\end{itemize}

\subsection{三、物理图像}\label{ux4e09ux7269ux7406ux56feux50cf-7}

圆筒的固有线电荷密度是 \(\lambda_1\)，在 \(S\) 中因纵向长度收缩变为
\(\gamma_v\lambda_1\)，并形成电流 \(I_1=\gamma_v\lambda_1v\)；轴线电荷
\(\lambda_2\)
静止，只贡献电场。圆筒内部高斯面不包围筒面电荷，安培环路也不包围筒面电流。

\subsection{四、解题入口}\label{ux56dbux89e3ux9898ux5165ux53e3-7}

先在 \(S\) 中用高斯面和安培环路分区求场。再把每一束电荷视作四维电流
\((c\lambda,j_z)\)，或等价地分别变换速度和线密度。

\subsection{五、解析}\label{ux4e94ux89e3ux6790-7}

\begin{figure}[H]
\centering
\includegraphics[width=0.706\linewidth,height=\textheight,keepaspectratio,alt={图 8-2　点 P 的横向电场、磁场及沿轴参考系变换。}]{../figures/solution-08.pdf}
\caption*{图 8-2　点 \(P\) 的横向电场、磁场及沿轴参考系变换。}
\end{figure}

图 8-2 标出 \(\bm E\) 的径向、\(\bm B\) 的环向以及参考系速度 \(u\)
的轴向，三者两两垂直。下面先按 \(d<R\)、\(d>R\)
判断高斯面和安培环路包围的源，再沿图示轴向作洛伦兹变换。

记

\[
\gamma_v=\frac1{\sqrt{1-v^2/c^2}},\qquad
\gamma_u=\frac1{\sqrt{1-u^2/c^2}}.
\]

题给圆筒随动截面上的固有线密度为

\[\frac{\lambda_1}{2\pi R}\cdot2\pi R=\lambda_1.\]

圆筒在 \(S\) 中沿轴运动，故 \(S\) 中测得的线密度是

\[\lambda_{1S}=\gamma_v\lambda_1.\]

\Needspace{3\baselineskip}
\textbf{1. \(S\) 系中的电场}

当 \(d>R\)，半径 \(d\)、长度 \(L\) 的高斯柱包围两部分电荷：

\[Q_{\rm enc}=(\gamma_v\lambda_1+\lambda_2)L.\]

由

\[E(2\pi dL)=\frac{Q_{\rm enc}}{\varepsilon_0},\]

得

\[
\boxed{
E(d>R)=\frac{\gamma_v\lambda_1+\lambda_2}{2\pi\varepsilon_0d}
}.
\]

正负号表示径向方向；若只写大小，应取绝对值。

当 \(d<R\)，高斯柱只包围轴线电荷：

\[
\boxed{
E(d<R)=\frac{\lambda_2}{2\pi\varepsilon_0d}
}.
\]

\Needspace{3\baselineskip}
\textbf{2. \(S\) 系中的磁场}

只有运动圆筒形成电流

\[I_1=\gamma_v\lambda_1v.\]

当 \(d>R\)，安培环路包围该电流：

\[B(2\pi d)=\mu_0I_1,\]

故

\[
\boxed{
B(d>R)=\frac{\mu_0\gamma_v\lambda_1v}{2\pi d}
}.
\]

当 \(d<R\)，环路位于筒面内部，不包围筒面电流：

\[\boxed{B(d<R)=0}.\]

\Needspace{3\baselineskip}
\textbf{3. \(S'\) 中的速度和线密度}

圆筒速度按共线速度变换：

\[
\boxed{v_1'=\frac{v-u}{1-uv/c^2}}.
\]

中轴导线在 \(S\) 中速度为零，因此

\[\boxed{v_2'=-u}.\]

圆筒电荷在自身静止系中的固有线密度就是 \(\lambda_1\)。在 \(S'\)
中圆筒的线密度为

\[
\boxed{\lambda_1'=\gamma_{v_1'}\lambda_1}.
\]

并且

\[
\gamma_{v_1'}=\gamma_u\gamma_v
\left(1-\frac{uv}{c^2}\right).
\]

轴线电荷在 \(S\) 中静止，故

\[
\boxed{\lambda_2'=\gamma_u\lambda_2}.
\]

\Needspace{3\baselineskip}
\textbf{4. 只有电场而无磁场}

在 \(d>R\) 区域，\(S'\) 中磁场为零等价于总轴向电流为零：

\[\lambda_1'v_1'+\lambda_2'v_2'=0.\]

分别代入：

\[
\gamma_{v_1'}\lambda_1
\frac{v-u}{1-uv/c^2}
-\gamma_u\lambda_2u=0.
\]

代入 \(\gamma_{v_1'}\) 并约去 \(\gamma_u\) 与公共因子：

\[\gamma_v\lambda_1(v-u)-\lambda_2u=0.\]

所以

\[
\boxed{
u_E=\frac{\gamma_v\lambda_1}{\gamma_v\lambda_1+\lambda_2}v
=\frac{\lambda_1}{\lambda_1+\lambda_2\sqrt{1-v^2/c^2}}v
}.
\]

该参考系必须满足

\[
\boxed{
\left|\frac{\lambda_1v}
{\lambda_1+\lambda_2\sqrt{1-v^2/c^2}}\right|<c
}.
\]

\Needspace{3\baselineskip}
\textbf{5. 只有磁场而无电场}

电场为零等价于 \(S'\) 中总线电荷为零：

\[\lambda_1'+\lambda_2'=0.\]

代入密度变换：

\[
\gamma_u\left[
\gamma_v\lambda_1\left(1-\frac{uv}{c^2}\right)+\lambda_2
\right]=0.
\]

因此

\[
\gamma_v\lambda_1+\lambda_2
-\gamma_v\lambda_1\frac{uv}{c^2}=0,
\]

\[
\boxed{
u_B=\frac{c^2}{v}
\frac{\gamma_v\lambda_1+\lambda_2}{\gamma_v\lambda_1}
=\frac{c^2}{v}
\frac{\lambda_1+\lambda_2\sqrt{1-v^2/c^2}}{\lambda_1}
}.
\]

存在条件为

\[
\boxed{
\left|\frac{c^2[\lambda_1+\lambda_2\sqrt{1-v^2/c^2}]}
{v\lambda_1}\right|<c
}.
\]

\Needspace{3\baselineskip}
\textbf{6. 能否同时没有电场和磁场}

若同一 \(S'\) 中两场同时为零，就要求 \(u_E=u_B\)。两式相乘得到

\[
u_Eu_B
=\frac{\lambda_1v}{\lambda_1+\lambda_2/\gamma_v}
\frac{c^2(\lambda_1+\lambda_2/\gamma_v)}{v\lambda_1}
=c^2.
\]

若两者相等，则

\[u^2=c^2,\]

即
\(|u|=c\)。有质量观察者的惯性系不能以光速运动，因此只要原位置的两场并非同时为零，任何合法惯性变换都不能将它们同时消去；原本无场的抵消情形须另列，见上面的退化分类：

\[\boxed{\text{不可能同时使 }\bm E'=0\text{ 且 }\bm B'=0}.\]

\Needspace{3\baselineskip}
\textbf{用场变换作独立核验}

在 \(d>R\) 区域，\(\bm E\) 径向、\(\bm B\) 环向，\(S'\)
沿轴向运动，三者两两垂直。适当选取正方向后

\[E'=\gamma_u(E-uB),\]

\[B'=\gamma_u\left(B-\frac{uE}{c^2}\right).\]

令 \(B'=0\) 得 \(u=c^2B/E\)，代入本题的 \(E,B\) 正好得到 \(u_E\)；令
\(E'=0\) 得 \(u=E/B\)，得到 \(u_B\)。两种方法相互验证了符号和 \(c^2\)
因子。

\Needspace{3\baselineskip}
\textbf{电磁场不变量的判据}

本题始终有 \(\bm E\cdot\bm B=0\)。另一不变量是

\[\mathcal I=E^2-c^2B^2.\]

若 \(\mathcal I>0\)，存在纯电场系而不存在纯磁场系；若
\(\mathcal I<0\)，结论相反；若 \(\mathcal I=0\)，所需变换速度达到
\(c\)，不存在普通惯性系能消去其中任一场。这与 \(u_Eu_B=c^2\)
的代数结论完全一致。

\Needspace{3\baselineskip}
\textbf{线密度变换的四维电流推导}

对圆筒电荷，\(S\) 系中的四维线流为

\[J^\mu=(c\gamma_v\lambda_1,0,0,\gamma_v\lambda_1v).\]

沿 \(z\) 方向以速度 \(u\) 变换：

\[
\begin{aligned}
c\lambda_1'
&=\gamma_u(c\gamma_v\lambda_1-u\gamma_v\lambda_1v/c)\\
&=c\gamma_u\gamma_v\lambda_1
\left(1-\frac{uv}{c^2}\right).
\end{aligned}
\]

因此

\[
\lambda_1'=\gamma_u\gamma_v\lambda_1
\left(1-\frac{uv}{c^2}\right)
=\gamma_{v_1'}\lambda_1.
\]

这个推导同时说明，电荷总量不变并不意味着不同参考系测得的线密度相同，因为同时测量的线段长度发生了洛伦兹收缩。

\Needspace{3\baselineskip}
\textbf{纯电场系与纯磁场系的条件互斥}

若 \(|u_E|<c\)，由 \(u_Eu_B=c^2\) 立即得到

\[|u_B|=\frac{c^2}{|u_E|}>c,\]

因此纯磁场候选系不具有物理意义。反之亦然。无需分别解复杂不等式，就能快速判断两类参考系最多存在一种。

\Needspace{3\baselineskip}
\textbf{分区场的边界跳变}

在 \(d=R\) 处，筒面电荷使法向电场满足

\[
E_{\rm out}-E_{\rm in}
=\frac{\gamma_v\lambda_1}{2\pi R\varepsilon_0}.
\]

表面电流使磁场满足

\[
B_{\rm out}-B_{\rm in}
=\frac{\mu_0\gamma_v\lambda_1v}{2\pi R}.
\]

两式恰等于正文内外场表达式之差，说明分区计算满足电磁边界条件。

\Needspace{3\baselineskip}
\textbf{四维线流的范数检查}

变换后圆筒电流为 \(j_1'=\lambda_1'v_1'\)，于是

\[
\left(c\lambda_1',j_1'\right)
=\gamma_{v_1'}\lambda_1(c,v_1').
\]

其闵可夫斯基范数

\[
c^2\lambda_1'^2-j_1'^2
=c^2\lambda_1^2
\]

保持不变，证明 \(\lambda_1\) 确实是圆筒静止系中的固有线密度。

\Needspace{3\baselineskip}
\textbf{面密度所在参考系必须先定清楚}

扫描题面把密度放在``在 \(S\) 中''的叙述里，而参考解又乘
\(\gamma_v\)，两者口径不一致。正文为了对应参考解，将 \(\lambda_1\)
明确为圆筒静止系密度。如果严格把给定值视为 \(S\)
系观测密度，则所有场公式中的 \(\gamma_v\lambda_1\) 要替换成
\(\lambda_1\)，并有

\[\lambda_1'=\gamma_u\lambda_1(1-uv/c^2),\qquad
u_E=\frac{\lambda_1v}{\lambda_1+\lambda_2},\quad
u_B=\frac{c^2(\lambda_1+\lambda_2)}{\lambda_1v}.\]

速度合成不变。不能既把 \(\lambda_1\) 当成 \(S\)
系给定量，又再乘收缩因子。

\Needspace{3\baselineskip}
\textbf{零源与退化情形}

记外区总密度 \(a=\gamma_v\lambda_1+\lambda_2\)、电流
\(j=\gamma_v\lambda_1v\)。若
\(a=j=0\)，外区原本就无场，任意亚光速参考系仍无场；这也包括
\(v=0\)、\(\lambda_2=-\lambda_1\ne0\)
的非零电荷抵消情形。原稿只列两密度均零并不完整。

若 \(j=0,a\ne0\)，\(S\) 已是纯电场系，不能使用除以 \(j\) 的 \(u_B\)
式；若 \(a=0,j\ne0\)，\(S\) 已是纯磁场系，不能使用除以 \(a\) 的 \(u_E\)
式。仅当 \(a,j\) 都非零时，才用两候选速度乘积 \(c^2\)。筒内 \(d<R\) 若
\(\lambda_2=0\)，则本来两场均零；若
\(\lambda_2\ne0\)，内区是纯电型，沿轴运动会产生环向磁场。

图 8-2 用横截面表示轴向速度出纸面、径向 \(E\) 和环向
\(B\)，原稿在横截面内画轴向速度向上不符合空间方向。

\Needspace{3\baselineskip}
\textbf{由总线流的不变量判断能否消去一种场}

将外部场对应的总线密度和电流记为 \(a,j\)。轴向变换直接给

\[a'=\gamma_u(a-uj/c^2),\qquad j'=\gamma_u(j-ua).\]

因此 \(c^2a'^2-j'^2=c^2a^2-j^2\)。若该量为正，可取 \(u=j/a\)
消去电流，得到纯电场；若为负，可取 \(u=c^2a/j\)
消去密度，得到纯磁场。两速度乘积为
\(c^2\)，非零场时不可能两者同时亚光速。

若不变量为零但 \((a,j)\ne(0,0)\)，候选速度的大小为
\(c\)，不存在物理惯性系消去其中一个场。若
\((a,j)=(0,0)\)，则外场本来就全为零，任意亚光速系仍为零。分母为零的特殊情形要先回到这两个线性方程，不能从候选速度分式推断不存在解。

筒内只包围静止轴线电荷，原系已是纯电；若
\(\lambda_2=0\)，筒内所有场在任何系都为零。这与外区的抵消条件分别判断。

\subsection{六、关键洞察}\label{ux516dux5173ux952eux6d1eux5bdf-7}

``纯电场系''和``纯磁场系''的候选速度互为 \(c^2/u\)。若一个候选速度小于
\(c\)，另一个必大于 \(c\)；这对应电磁场不变量 \(E^2-c^2B^2\)
的正负分类。

相对论连续电荷题最稳妥的对象是四维电流，而不是孤立地背``密度乘某个
\(\gamma\)''。密度与电流会在洛伦兹变换中彼此混合。

\textbf{先区分三种线密度。} \(\lambda_1\)
是圆筒电荷在自身静止系中的固有线密度；\(S\)
系看到圆筒沿轴运动，纵向长度收缩，故密度为 \(\gamma_v\lambda_1\)；\(S'\)
系还要使用圆筒相对 \(S'\) 的速度 \(v_1'\)，密度变为
\(\gamma_{v_1'}\lambda_1\)。同一个符号若在三个参考系间不加说明，会使后续所有电场和电流条件混乱。

\textbf{内外区域由包围源决定。} 对
\(d>R\)，高斯柱包围筒面电荷和轴线电荷，安培环路也包围圆筒表面电流；对
\(d<R\)，高斯柱只包围轴线电荷，安培环路不包围位于 \(R\)
的表面电流。这里不需要猜场是否``穿过''圆筒，只需严格看积分曲面包围了哪些源。

\textbf{纯场参考系的物理条件。} 在轴对称外区，\(B'=0\) 等价于 \(S'\)
系总轴向电流为零；\(E'=0\)
等价于总线电荷密度为零。一个是电流四矢量的空间分量消失，一个是时间分量消失。两者是否存在由四维电流是类时还是类空决定，对应场不变量
\(E^2-c^2B^2\) 的正负。

\textbf{不变量给出最后判决。} 若某系中可令 \(B'=0\)，则
\(E^2-c^2B^2>0\)；若可令
\(E'=0\)，则该不变量小于零；若等于零，场为零型，只有在以光速极限追赶时才可能逼近纯场，普通惯性系不能同时消去两者。因此式
\(u_Eu_B=c^2\) 不是偶然代数，而是场分类的直接反映。

\textbf{四维电流把两种条件统一起来。}
对沿轴分布的一束线电荷，在某参考系中记线密度为 \(\lambda\)、电流为
\(I=\lambda w\)，可把它们组成二维四维电流

\[
J^\mu=(c\lambda,I).
\]

沿轴以速度 \(u\) 作洛伦兹变换：

\[
\begin{aligned}
c\lambda'&=\gamma_u
\left(c\lambda-\frac{u}{c}I\right),\\
I'&=\gamma_u(I-u\lambda).
\end{aligned}
\]

即

\[
\lambda'=\gamma_u\left(\lambda-\frac{uI}{c^2}\right),
\qquad
I'=\gamma_u(I-u\lambda). \tag{A}
\]

密度和电流不是分别变换的两个无关标量，而是同一四维对象的两个分量。

\textbf{把式 (A) 用于两束电荷。} 在 \(S\)
系中，圆筒电荷的观测线密度与电流为

\[
\lambda_{1S}=\gamma_v\lambda_1,
\qquad
I_{1S}=\gamma_v\lambda_1v,
\]

轴线电荷为

\[
\lambda_{2S}=\lambda_2,
\qquad I_{2S}=0.
\]

先求总量再变换：

\[
\begin{aligned}
\lambda'_{\rm tot}
&=\gamma_u\left[\gamma_v\lambda_1+\lambda_2
-\frac{u}{c^2}\gamma_v\lambda_1v\right],\\
I'_{\rm tot}
&=\gamma_u\left[\gamma_v\lambda_1v
-u(\gamma_v\lambda_1+\lambda_2)\right].
\end{aligned} \tag{B}
\]

令 \(I'_{\rm tot}=0\)，得到无磁场候选速度

\[
u_E
=\frac{\gamma_v\lambda_1v}
{\gamma_v\lambda_1+\lambda_2}
=\frac{\lambda_1v}
{\lambda_1+\lambda_2/\gamma_v}.
\]

令 \(\lambda'_{\rm tot}=0\)，得到无电场候选速度

\[
u_B
=\frac{c^2(\gamma_v\lambda_1+\lambda_2)}
{\gamma_v\lambda_1v}
=\frac{c^2}{v}
\frac{\lambda_1+\lambda_2/\gamma_v}{\lambda_1}.
\]

二者相乘立即得到 \(u_Eu_B=c^2\)。这条路线比逐束先求 \(v_i'\)、再求
\(\gamma_{v_i'}\) 更短，但两种方法应得到完全相同的结果。

\textbf{何时选哪条路线。} 若题目直接给某一参考系中的总密度和总电流，式
(A)
最简；若给的是每束电荷的固有密度，则逐束速度变换更不易混淆。课堂上可先用逐束表格建立物理图像，再用总四维电流压缩计算，形成``可理解''和``可高效''两种层次。

\subsection{七、易错点}\label{ux4e03ux6613ux9519ux70b9-7}

\begin{itemize}
\tightlist
\item
  把圆筒固有线密度直接当成 \(S\) 系线密度，漏掉 \(\gamma_v\)。
\item
  在 \(d<R\) 时仍把筒面电荷和电流计入高斯面、安培环路。
\item
  写 \(\lambda_1'=\gamma_u\lambda_1\)，没有使用圆筒相对 \(S'\) 的速度
  \(v_1'\)。
\item
  共线速度变换写成简单的 \(v-u\)。
\item
  求出 \(u\) 后不检验 \(|u|<c\)。
\item
  只变换速度而不变换线密度，导致总电流条件错误。
\item
  场变换中把 \(uE/c^2\) 漏写成 \(uE\)。
\item
  不区分 \(d<R\) 与 \(d>R\)，把同一组场公式用于全部空间。
\item
  把固有线密度的符号与速度方向混在一起，错误地把负电荷运动方向再反号一次。
\item
  得到 \(|u|=c\) 后仍把它当作合法惯性参考系。
\end{itemize}

\subsection{八、课堂讲解建议}\label{ux516bux8bfeux5802ux8bb2ux89e3ux5efaux8bae-7}

先用二维表列出两束电荷在 \(S\) 中的``线密度---速度---电流''，再逐行变到
\(S'\)。第3、4问分别解释为总电流为零和总电荷为零，物理意义会比直接套场变换更清楚。

完成电荷---电流法后，再用 \(E',B'\)
变换各做一次核验。两条路线得到同一速度，是学生建立相对论电磁统一观的好机会。

\textbf{课堂推导顺序。} 第一张表写 \(S\) 系：圆筒固有密度
\(\lambda_1\)、观测密度 \(\gamma_v\lambda_1\)、速度 \(v\)、电流
\(\gamma_v\lambda_1v\)；轴线电荷密度 \(\lambda_2\)、速度 0、电流
0。第二张表写 \(S'\) 系，再分别令总电流或总密度为零。

\textbf{为什么要有 \(\gamma_v\)。} 取圆筒上一段随动电荷，总电量不变，而
\(S\) 系同时测得的纵向长度因洛伦兹收缩变短，所以线密度增大为
\(\gamma_v\lambda_1\)。如果忽略这一点，后续纯场参考系速度会少掉
\(\sqrt{1-v^2/c^2}\)。

\textbf{存在性判断。} 每求出一个候选 \(u\)，立即检查 \(|u|<c\)。然后用
\(u_Eu_B=c^2\) 一步判断：两种纯场系不可能同时为亚光速。最后再用
\(E^2-c^2B^2\) 解释这种互斥的几何本质。

\textbf{两张表的具体格式。} 第一张表的列为``电荷组、固有密度、\(S\)
系速度、\(S\) 系密度、\(S\) 系电流''；第二张表把速度列替换为
\(v_i'\)，密度写成
\(\gamma_{v_i'}\lambda_i\)。表格最后一行求和，分别得到
\(\lambda'_{\rm tot}\) 与
\(I'_{\rm tot}\)。学生一旦把每束电荷独立变换，纯电和纯磁条件就只是令表格总和为零。

\textbf{速度合成不可跳步。} 圆筒在 \(S'\) 中速度

\[
v_1'=\frac{v-u}{1-uv/c^2},
\]

轴线电荷在 \(S'\) 中速度 \(v_2'=-u\)。建议先只判断三个极限：\(u=0\) 时
\(v_1'=v\)；\(u=v\) 时圆筒静止；\(v=0\) 时两束都以 \(-u\)
运动。三个极限都正确后再代入密度与电流，能有效避免分母符号错误。

\textbf{用场变换交叉核验。} 对互相垂直的径向电场和环向磁场，沿轴以 \(u\)
变换时

\[
E'_{\perp}=\gamma_u(E-uB),
\qquad
B'_{\perp}=\gamma_u\left(B-\frac{uE}{c^2}\right).
\]

令 \(B'=0\) 得纯电场系候选速度 \(u_E=c^2B/E\)；令 \(E'=0\)
得纯磁场系候选速度 \(u_B=E/B\)。将 \(S\) 系的 \(E,B\)
代入，应与四维电流法一致。这是本题最重要的双路线核验。

\textbf{符号与存在性板书。} 不要在最初就把 \(\lambda_1,\lambda_2\)
假定为正。保留代数符号求出候选 \(u\)
后，再讨论分母是否为零、方向正负以及
\(|u|<c\)。异号电荷可能使总电荷或总电流抵消，也可能使候选速度超光速；存在性是答案的一部分。

\subsection{九、变式训练}\label{ux4e5dux53d8ux5f0fux8badux7ec3-7}

\begin{enumerate}
\def\labelenumi{\arabic{enumi}.}
\tightlist
\item
  若轴线电荷也以速度 \(w\) 运动，重做纯电场系的条件。
\item
  用场变换公式直接推出 \(u_E=c^2B/E\) 与
  \(u_B=E/B\)，并与电荷---电流方法核对。
\item
  取一组具体的异号
  \(\lambda_1,\lambda_2\)，判断存在的是纯电场系还是纯磁场系。
\end{enumerate}

\Needspace{3\baselineskip}
\textbf{变式解析要点}

\begin{enumerate}
\def\labelenumi{\arabic{enumi}.}
\tightlist
\item
  轴线电荷的四维电流不再是 \((c\lambda_2,0)\)，而要使用固有密度与速度
  \(w\)
  构造。先分别求两束在目标系中的密度和速度，再令总电流为零；纯磁场系才令总密度为零。
\item
  场变换法只需外区 \(E\perp B\perp\hat{\bm x}\)
  的两式；求得候选速度后仍须检查方向约定和
  \(|u|<c\)，不能因为公式短就省略存在性。
\item
  选定数值后先计算
  \(E^2-c^2B^2\)。正值预示纯电场系，负值预示纯磁场系，零值预示两者都不能在亚光速惯性系中完全消去；再用显式
  \(u\) 验证。
\end{enumerate}
\end{document}
